\chapter{知识的更新和维护}
\label{chap:ket-update-maintenance}

花卉识别系统在春季采集的照片上表现正常，到了雨季却频繁出错。此时，用户仍然要求识别原来的物种，系统却可能需要重新学习：照片背景变了，物种出现的比例变了，或者新的观测遗漏了原先用于区分类别的信息。另一个植物资料问答系统也可能需要更新，但原因不是照片分布变化，而是某研究机构的负责人发生了变更。若有人撤回一条曾经参与训练的记录，仅把答案改成拒答又不能满足消除训练影响的要求。这些问题都发生在已有系统中，却不能靠同一种操作、同一个分数来判定是否解决。

本章讨论新证据或修改要求到来以后，怎样改变系统，同时维持仍需提供的能力。第~\ref{chap:ket-reuse-selection}章已经讨论复用价值与候选知识的选择；这里固定已有系统及其维护责任，依次确定更新目标、准备证据、实施操作、验证结果并维护连续版本。这个过程允许返回：验证发现历史回答被误改时，需要重新界定修改范围；当前标签不足以支持更新时，需要补充证据，而不是一味增加训练次数。对尚未见过的环境，也可以利用已知环境准备下一版本，但不能把未来目标数据偷偷纳入训练或选模。

以下沿用来源和目标联合分布$\mathcal P_S,\mathcal P_T$、相应风险$R_S,R_T$以及经验风险$\hat R$。模型写作$f_{\bm{\theta}}$，输入为$\bm{x}$，表示为$\bm{z}$；学习阶段或迭代使用上标$^{(t)}$，任务$\mathcal T_k$使用下标$k$。环境变化下可学习性的假设和不可识别边界见第~\ref{chap:transfer-learning-theory}章。本章在这些边界内展开可实施的更新。若物种集合、输出目标或成功标准扩展，则由第~\ref{sec:ket-transfer-output-semantics}节确定新任务，再调用这里的保留与维护机制；已有模型的整合和多任务共同训练分别见第~\ref{chap:ket-integration}章与第~\ref{chap:ket-sharing}章。

更新路线不能只按算法名称选择。一次请求应同时落在四个坐标上：目标是环境适应、事实修正、持续吸收还是训练影响消除；载体是数据与外部记录、表示与参数、模块与路由还是运行状态；信息权限是可读旧数据、可读参数、只查询输出或仅有当前流；时间协议是离线批量、阶段序列、严格在线还是带有效区间的版本请求。表~\ref{tab:ket-update-route-matrix}给出主要交叉，单次维护可以占多格，不能把行列强改成互斥类别。

\begin{table*}[htbp]
\small
\begin{tabularx}{\linewidth}{@{}p{23mm}XXXX@{}}
\toprule
更新目标 & 优先检查的载体 & 关键权限 & 时间协议 & 路线选择问题 \\
\midrule
环境适应 & 预处理、表示、参数与统计状态 & 当前输入；标签或源接口按协议声明 & 离线批量、测试时或严格在线 & 变化是否真的破坏任务规律，代理信号能否支持更新 \\
事实修正 & 外部记录、索引、缓存或参数行为 & 有效请求、来源、依赖与历史查询 & 生效区间和连续版本 & 改哪里能同时满足新事实、历史语义与局部性 \\
持续吸收 & 样本记忆、参数、模块和优化状态 & 历史访问、容量与反馈时机 & 阶段序列、环境回归或无边界流 & 怎样兼顾当前吸收、旧能力与后续可塑性 \\
训练影响消除 & 原记录及派生物、参数与恢复快照 & 撤回集合、剩余数据和规定学习过程 & 请求到达后持续有效 & 外部删除与参数影响消除分别需要什么证据 \\
\bottomrule
\end{tabularx}
\caption{更新目标、知识载体、信息权限与时间协议的路线矩阵。各行是主要责任而非互斥分类；例如在线事实修正可能同时修改外部记录、缓存和参数，并须满足持续维护要求。}
\label{tab:ket-update-route-matrix}
\end{table*}

\section{知识更新目标的确定}
\label{sec:ket-update-goals}

观察到变化，并不意味着必须改变模型；决定改变模型，也不意味着每个旧答案都应受到保护。更新前需要回答三个相互约束的问题：失效究竟来自什么，哪些内容允许改变，以及哪些能力和资源限制必须保持。只有把这些问题分开，才可能判断一次更新是在修复问题，还是改变了原本要解决的问题。

\subsection{更新需求的辨认}
\label{sec:ket-update-needs}

设原系统和当前系统都把照片映射到同一组花卉物种。若输入$\bm{x}$的分布改变，而给定照片后的标签规律不变，即$p_T(y\mid\bm{x})=p_S(y\mid\bm{x})$，便得到第~\ref{sec:distribution-shift}节的\term{协变量偏移}（covariate shift）：原有判断依据仍然成立，只是某些照片更常见了。若保持的是$p_T(\bm{x}\mid y)=p_S(\bm{x}\mid y)$，而$p_T(y)$改变，则是\term{标签偏移}（label shift）：每类照片内部的统计特征不变，变化的是各类在总体中的份额。这两个条件一般不能互换，标签比例改变通常也会改变后验$p(y\mid\bm{x})$。

当$p_T(y\mid\bm{x})$改变时，应继续检查它能否由标签偏移解释。若类别内部的观测规律也变了，就不能再只校正类别比例，需要更一般的证据。即使物种的生物学定义没有改变，输入中未记录的因素也可能使条件分布改变。例如，系统只保存经过压缩的外观特征，两种近似物种在这些特征上难以区分：仅改变它们的出现比例可能符合标签偏移；相机处理流程丢失细微纹理，则还会改变类内观测规律。两者都可能改变相同记录特征所对应的物种概率，却需要不同的修正依据。这里的条件规律是相对于\emph{实际记录的输入}定义的，不是对完整现实对象的断言。若变化来自标注标准重新划分物种，则还要核查输出语义是否已变，不能直接作为旧任务偏移处理。

统计关系变化与\term{观测空间变化}（change of observation space）可以同时出现。后者表示系统取得同类任务信息的接口改变。本章把概述中的传感器更换具体化为从RGB照片改用深度图，便于检验新观测能否保留花色信息；它不是多光谱传感器的同一组测量。即使把两种输入都展平成相同长度的向量，颜色强度与距离数值也不是同一种含义；反过来，同一照片的无损编码长度改变，并不必然改变可用于判断的信息。跨观测更新需要说明同一植株、局部器官或类别语义怎样对应，以及颜色等旧信息能否从新传感器中恢复。

可以用三次部署变化检验这些判断。若仅把拍摄地点从温室换到露天，输入和物种含义保持不变，但背景、遮挡和类别比例可能同时改变，不能只因“换了背景”就宣称满足协变量偏移；需要检验或明确采用的近似。若未记录的生长阶段改变相同外观特征下的类别概率，输入与目标的名称虽未改变，条件预测规律已经改变。若改用深度相机，识别目标仍是原物种，但输入含义和可见信息改变，需要建立跨观测对应。仿真到现实若只改变照片形成条件，可采用相同分析；动力学、奖励和交互后果还需专门建模，这属于第~\ref{part:reinforcement-learning}部分的内容范围。

变化检测只能产生调查线索。假设识别器仅依赖花瓣形态，背景从灰色变成绿色，而形态与标签的关系及观测质量保持不变，输入分布虽有变化，原规则仍可能准确，不更新就是合理候选。相反，准确率下降也不必然来自模型知识过时：标签录入错误、深度值单位变化、照片解码故障以及评价类别权重变化，都可能造成下降。诊断应先核对数据、接口和评价口径，再决定采用\term{领域自适应}（亦称领域适应，domain adaptation），即利用当前环境可取得的证据调整系统，使原任务在该环境继续有效。

\subsection{修改内容与适用范围的确定}
\label{sec:ket-update-scope}

环境适应通常由一批数据触发，定向修改则可能只由一条请求触发。为了知道该改变什么，可以把请求记为
\begin{equation}
e=(s,r,o_{\mathrm{old}},o_{\mathrm{new}},c,[\tau_0,\tau_1),a),
\label{eq:ket-update-request}
\end{equation}
其中$s$是主体身份，$r$是关系，$o_{\mathrm{old}},o_{\mathrm{new}}$是当前与期望记录，$c$表示其他适用条件，$[\tau_0,\tau_1)$是有效时间，$a$是请求依据。时间$\tau$指事实适用的日历时间，不是参数迭代编号$t$。这份记录既限制修改范围，也提供验证答案，不能缩成“遇到某个词就换一个词”。

考虑贯穿本章的教学构造：植物研究机构甲所的负责人原为人员甲，自2026年7月1日起改为人员乙。为避免名字混淆，机构和人员均具有独立身份标识。假设制度明确规定负责人担任项目签批人，且对外联系应转至负责人办公室，那么当前负责人、当前签批人和当前联系办公室都应随生效时间改变。查询“2026年6月15日由谁负责”仍应回答人员甲；查询“2026年7月15日由谁负责”应回答人员乙。人员甲负责另一家同名机构的记录、甲所的成立年份和已完成历史审批的签字人，都不因这次任命自动改变。

图~\ref{fig:ket-update-time-scope}把时间和依赖放在同一张图中。依赖边不是语言相似度，而是已声明的推理规则：负责人发生变化之所以影响签批人，是因为例子给出了相应制度。现实系统无法仅凭两个句子包含相同实体就判断传播范围；没有证据支持的关联结论应保持待决，不能自行补造。

\begin{figure*}[htbp]
\centering
% Teaching construction: appointment effective on 2026-07-01.
% Solid arrows are only the declared dependencies, not inferred associations.
\begin{tikzpicture}[
  x=1mm,y=1mm,font=\figurefont\small,
  box/.style={rectangle,draw=BookInk,fill=BookPaper,line width=0.7pt,
    align=center,inner xsep=2mm,inner ysep=2mm,minimum height=12mm},
  flow/.style={draw=BookInk,line width=0.8pt,-{Latex[length=2mm]},
    rounded corners=1mm}
]
  \node[anchor=west,text=BookInk] at (8,84) {2026年7月1日前};
  \node[anchor=west,text=BookInk] at (84,84) {2026年7月1日起};
  \draw[BookBlue,line width=2pt] (8,77) -- (80,77);
  \draw[BookTeal,line width=2pt] (80,77) -- (156,77);
  \draw[BookInk,line width=0.7pt] (80,74) -- (80,80);

  \node[box,text width=60mm,fill=FigureInput!10] (old) at (42,60)
    {旧事实：甲所负责人为甲\\有效至生效点之前};
  \node[box,text width=60mm,fill=FigureOutput!10] (new) at (122,60)
    {新事实：甲所负责人为乙\\自生效点开始有效};
  \node[box,text width=60mm] (history) at (42,34)
    {历史查询：2026年6月15日\\负责人仍为甲};
  \node[box,text width=26mm] (approve) at (101,34)
    {当前签批人\\变为乙};
  \node[box,text width=26mm] (contact) at (143,34)
    {当前联系对象\\乙的办公室};
  \draw[flow] (old.south) -- (history.north);
  \draw[flow] (new.south) -- (122,47) -| (approve.north);
  \draw[flow] (122,47) -| (contact.north);
  \node[box,text width=60mm,draw=BookMuted] at (42,10)
    {无关事实：甲所成立年份\\保持不变};
  \node[align=center,text width=70mm,text=BookMuted,
    font=\figurefont\footnotesize] at (122,10)
    {实线：给定的推理依赖\\无关事实不接依赖边};
\end{tikzpicture}

\caption{修改内容的时间与依赖范围。教学示例以2026年7月1日为任命生效点，时间条区分旧事实与新事实的有效区间；实线箭头表示已经给定的制度依赖。当前签批和联系对象需要更新，历史查询仍使用其时间范围内的旧事实，无关成立年份不变。图中仅列已知依赖，不表示系统已经识别全部事实依赖。}
\label{fig:ket-update-time-scope}
\end{figure*}

\term{事实编辑}（factual editing）要求改正指定事实及其相关使用；\term{行为修正}（behavior correction）要求改变特定条件下的响应，例如缺少有效任命来源时明确说明无法确认，而不是猜测现任负责人。两者都要规定允许变化的语境。若行为要求改变了系统的基本成功标准，则任务定义由第~\ref{sec:ket-transfer-output-semantics}节处理；本章负责在要求已经明确以后实施修正。

\term{机器遗忘}（machine unlearning）具有另一种目标：消除指定数据对训练结果的影响，而不只是改变某个问答输出。撤回请求应标明记录身份、所属训练版本及派生副本，例如原照片、增强样本、特征缓存和由它生成的训练标签。需要撤回一条含有人员姓名的训练记录，不等于禁止模型在所有语境提到该姓名；其余合法数据仍可能支持相同事实。反过来，屏蔽姓名也不证明该记录对参数和其他输出的影响已经消除。严格的请求范围和删除保证见第~\ref{sec:privacy-unlearning}节。

\subsection{保留要求与资源边界的确定}
\label{sec:ket-update-requirements}

“更新后更好”至少包含三种不同判断。\term{能力保留}（retention）考察仍需服务的旧任务是否保持；\term{可塑性}（plasticity）考察系统以后还能否用新证据学会需要的内容；\term{知识积累}（knowledge accumulation）则进一步要求历史经验降低后续学习代价，或者新经验改善仍有效的旧能力。旧任务分数不下降，但新标签再多也学不进去，说明保留不能替代可塑性；当前雨季识别变好但温室识别变差，则可能是适应伴随遗忘。即使两者都没有退化，也还需要与没有相关历史经验的参照比较，才能声称出现正向积累。

保留并非复制全部旧响应。对于询问2026年7月15日甲所负责人的问题，旧答案甲正是应当被修正的对象，不能同时纳入“必须严格保持”的集合。局部性限制的是不必要的外溢，而不是阻止必要的依赖传播；撤回目标也优先于保留受撤回数据支持的历史行为。验收清单因而应区分继续服务的任务、允许改变的输入条件、必须吸收的新信息以及禁止重新引入的训练记录。

表~\ref{tab:ket-update-goals}把几类目标对应到具体对象。一个请求可以同时涉及环境适应和持续吸收，也可以同时涉及事实修正与记录撤回；此时不是挑选一行代表整次更新，而是分别准备每项要求的证据。一次编辑通过，只说明这个请求及其测试范围内的结果；批量编辑要检查请求间冲突，连续更新还要检查前后操作覆盖、记忆增长和恢复能力。

\begin{table*}[htbp]
\small
\begin{tabularx}{\linewidth}{@{}p{16mm}XXXXX@{}}
\toprule
更新目标 & 期望变化 & 仍需保持的对象 & 允许的信息 & 比较参照 & 主要误判 \\
\midrule
环境适应 & 当前环境原任务风险下降 & 仍服务的原环境与类别语义 & 声明的目标输入、标签及源信息 & 原版本与不更新基线 & 把更自信当成更准确 \\
事实修正 & 指定条件下使用新事实 & 历史事实、无关关系与任务 & 有效请求、来源、时间和依赖 & 依有效条件制定的正确答案 & 只测原提示，忽略历史和关联 \\
持续吸收 & 学会当前内容并能继续学习 & 有效旧能力及可用学习容量 & 规定的经验流与记忆 & 等样本、等预算的阶段对照 & 只看遗忘，不测后续学习 \\
训练记录撤回 & 消除指定记录的训练影响 & 剩余任务及其他有效修改 & 剩余数据与允许的训练状态 & 剩余数据上的指定学习过程 & 把拒答或记录删除当成影响消除 \\
\bottomrule
\end{tabularx}
\caption{更新目标、保留要求与验收对象。“允许的信息”须在操作前确定，表中的可用依据不自动满足授权或隐私要求。多目标请求需要同时满足相应行的验收要求。}
\label{tab:ket-update-goals}
\end{table*}

资源边界同样应事先确定。固定总存储、固定更新次数、限制单次停机时间，会分别限制保存多少证据、重复使用多少经验以及是否允许完整重训练。旧模型、优化器状态、生成器、编辑器、外部索引和验证副本都要计入相应预算。只比较在线一次梯度更新的时间，却忽略编辑器预训练或历史模型保存成本，会改变比较问题；在结果出来以后才排除一种方法必需的成本，也不能形成一致的结论。

\section{知识更新依据的准备}
\label{sec:ket-update-evidence}

目标已经明确以后，仍要判断什么材料能够支持它。证据准备不只是在磁盘上找到一个数据集，还包括确认材料来自哪个版本、适用于什么时间、允许怎样使用，以及它能否辨认当前失配。保存依据与使用依据是两件事：本节确定保存什么，后文再说明这些材料怎样进入取样、目标和参数更新。

\subsection{可用信息与访问条件的确认}
\label{sec:ket-update-access}

当前样本、历史样本、旧预测、旧表示、模型参数、外部记录和修改请求提供不同证据。源数据不可访问，仍可能有可读取的源模型；只能查询旧预测，不意味着能够计算源模型中间表示。没有源数据、没有目标标签和没有目标输入也是三个不同限制：前两种设定还可能利用当前无标签照片，第三种则只能依据已知环境准备更新。每份材料都应带有来源、生成版本、适用时间和允许用途，避免把不同模型或不同有效条件下的输出混为同一参照。

SHOT提供了“源数据不可用，但源模型可读”的例子。它保留源分类器，调整目标侧特征提取器，使新表示能被旧分类器使用；无标签目标数据通过预测信息最大化和由目标特征原型得到的伪标签提供信号。这里冻结分类器不等于没有访问内部参数，因为更新特征提取器需要通过分类器反向传播。原型和信息目标也不能凭空确定任意目标标签语义\scite{ket-liang2020-shot}

DINE处理的是只能查询源预测的条件。它先查询目标样本的源预测，将这些响应传给可训练的目标模型，再利用目标数据结构继续调整。其原文使用自适应标签平滑处理有限的概率输出，并在蒸馏时加入插值一致性和互信息约束；随后不再依赖源输入数据做目标侧微调。学生与源模型不必同构，但源概率所包含的软信息不能由一个类别编号完整恢复\scite{ket-liang2022-dine}

把同一个分类器放到三个接口后，可以更清楚地看出操作边界。读取参数且具有模型结构、预处理和可微实现时，通常能计算中间特征及输入梯度；仅查询概率时，可以在获准提交的输入上记录软预测，却不能直接取得内部梯度，有限差分查询也不等于原生反向传播；仅查询类别时，能够记录决策结果，不能知道第二候选类别的概率。三个接口都只能对实际持有且允许提交的旧样本计算旧样本预测。查询预算、输入权限和模型版本不允许时，“有预测接口”不能补齐这些条件。

\begin{table*}[htbp]
\small
\begin{tabularx}{\linewidth}{@{}p{26mm}cccccX@{}}
\toprule
操作 & 历史原样本 & 模型参数 & 中间表示 & 源响应 & 当前标签 & 额外依赖与成本 \\
\midrule
原样本监督回放 & 必 & 学 & — & — & — & 需旧标签；保存样本并反复前后向计算 \\
历史特征回放 & — & 学 & 必 & — & — & 需旧特征及版本；只能训练特征之后的部分，除非另有原输入 \\
两域表示统计对齐 & 条 & 学 & 必 & — & — & 需两域输入或兼容的统计；编码、统计估计与批量成本 \\
重要性参数约束 & — & 必 & — & — & — & 需固定参考参数和先前估计的重要性；各需存储 \\
概率行为模仿 & — & 学 & — & 必 & — & 需输入及对应软预测；可在线查询或读取兼容缓存，另计学生训练 \\
类别行为模仿 & — & 学 & — & 必 & — & 需输入及对应硬标签，可查询或读取缓存；不足以直接匹配完整概率 \\
当前监督更新 & — & 学 & — & — & 必 & 需当前输入；标签取得和质量核对有独立成本 \\
\bottomrule
\end{tabularx}
\caption{信息访问条件与可实施操作。“必”为直接依赖，“—”为无此必要条件，“学”表示要能学习被更新模型的参数而不必读取源模型，“条”表示历史原样本可被兼容表示或统计替代。源响应可由在线查询或已保存的兼容缓存提供；使用缓存时，学习阶段无需再访问教师接口，但仍需对应输入。表示对齐还须明确表示来源。“可实施”不等于“有效”。}
\label{tab:ket-update-access}
\end{table*}

表~\ref{tab:ket-update-access}中，参数约束所需的重要性必须先有估计依据，不能由旧参数本身自动得到。概率匹配所需的源响应也应记录查询日期、服务版本、标签顺序、温度及能确认的预处理设置；未知配置明确留空。若服务途中换版，先后响应就未必来自同一教师。访问限制还不等同于隐私保证，模型或输出所泄露的信息需按第~\ref{chap:trust-privacy}章另外分析。

\subsection{历史证据的选择与保存}
\label{sec:ket-update-memory}

原始照片保留的信息最多，可以用新的编码器重新计算特征，但存储、授权与敏感内容的负担也较大。历史特征更紧凑，却只保留旧编码器留下的部分信息；如果新版本需要恢复被旧编码器丢弃的纹理，就不能靠这些特征补回。历史预测保存旧决策及类间相似性，但不是正确标签的替代物。生成模型可以按需产生旧样本，却把覆盖问题转移到生成质量和教师输出上，生成器权重及生成时间也不能不计成本。

在有限记忆中，代表性并不只有一种定义。随机保留近似反映经验流的出现频率，按特征中心保存强调常见形态，稀有类别配额保护低频但必须服务的场景，按预期干扰选择则强调容易在下一次更新中受损的样本。不确定样本可能暴露边界，也可能只是观测错误或异常值，因此不应把“越不确定越值得保存”当成通则。这些规则有不同覆盖目标，选定后还需要检查实际记忆遗漏了什么。

以容量为三张照片的教学记忆说明在线替换。前面依次看到晴天类别甲照片$a_1$、晴天类别甲照片$a_2$、阴影类别乙照片$b_1$，记忆先装满。第四张是常见晴天甲$a_3$，第五张是罕见遮挡乙$b_2$。若采用容量为$m$的蓄水池抽样，前$m$条记录依次入库；第$i>m$条到来时，从$\{1,\ldots,i\}$均匀抽取整数$j$，仅在$j\leq m$时用新记录替换第$j$个槽位。因此新记录以$m/i$的概率进入，进入时均匀替换一个槽位。设本次示例在$i=4$时抽到$j=2$、在$i=5$时抽到$j=5$，则第四步替换$a_2$，第五步不入库，最终为$\{a_1,a_3,b_1\}$。这是一条明确的随机实现，不表示算法必然丢弃遮挡样本；每个历史样本被保留的边际概率均为$3/5$。

若改为“保留晴天甲、阴影乙、遮挡乙三个已知情形”的覆盖规则，第四步可用$a_3$替换同情形中的$a_2$，第五步再用$b_2$替换冗余晴天甲，得到$\{a_1,b_1,b_2\}$。它保护了稀有情形，但记忆比例不再代表真实类别比例。若规则改为保留预估损失增长最大的照片，还需用当时的模型、候选更新和可访问数据估计干扰；第五步之后不能回看已经丢弃的$a_2$并宣称在线规则本来会选它。保存规则决定证据，训练权重决定这些证据的作用，两者需要分别记录。

iCaRL在类别增量设定中保存每类有序代表样本，用样本均值逼近类别特征均值，随着类别增多缩减每类名额。设当前可见类别$k$的归一化特征均值为$\bm{\mu}_k$，已选$r-1$个样本，则第$r$个选择使
\begin{equation}
\left\lVert\bm{\mu}_k-
\frac{1}{r}\left(\phi(\bm{x})+
\sum_{j=1}^{r-1}\phi(\bm{x}_{k,j})\right)\right\rVert_2
\label{eq:ket-update-herding}
\end{equation}
最小的未选样本；这里省略均值再归一化的实现细节，以突出逐项逼近的选择逻辑。新类别的完整当前批量可用于估计均值，旧类别却只能保留已有列表的前缀，不能假定以后还能访问全部旧样本。iCaRL还使用这些原样本重新计算表示并进行最近类均值预测，类别扩展的任务条件见第~\ref{sec:ket-transfer-output-semantics}节\scite{rebuffi2017}

固定总容量$m$时，$K$类平均只能分到约$m/K$个样本；固定每类$m_0$个时，总样本数是$Km_0$，并非同一个资源设定。标签、时间、模型版本和旧预测也是记忆的一部分。保存$\phi_{\bm{\theta}^{(t)}}(\bm{x})$时应记录$t$，因为以后计算的$\phi_{\bm{\theta}^{(t+1)}}(\bm{x})$可能已使用不同坐标关系；直接混合两代特征会让距离和均值失去一致含义。没有原样本时，只能冻结相应接口、估计兼容映射或承认可更新范围受限。

\subsection{当前证据与修改请求的核对}
\label{sec:ket-update-current-evidence}

少量当前标签能够回答无标签统计无法回答的问题。例如，两类花卉在目标特征空间中形成两个簇，未标注数据可以显示簇的位置，却不能保证哪个簇对应哪种物种。给每个簇补入可靠标签可以消除一种语义排列歧义，但未覆盖区域仍可能存在其他规律。MME利用已有源标签、少量目标标签和目标无标签样本，采用原型式分类器；有标签损失维持分类含义，无标签部分交替让分类器增大预测熵、让特征提取器减小预测熵，使原型与目标表示相互调整。它是在\emph{利用既有少量监督}，而不是决定应该标哪张照片\scite{ket-saito2019-mme}

CLUE则直接处理新增监督的选择：先计算目标特征和预测不确定性，再做不确定性加权聚类，选取接近各簇中心的实例请求标签，以兼顾信息量与覆盖。标签返回后再更新模型，因此应同时计入标注数、聚类和特征提取成本。若旧模型对某种系统错误过度自信，单靠不确定性权重仍可能遗漏它；覆盖性抽查和任务反馈提供不同证据。标注与主动查询的一般机制属于第~\ref{part:efficient-knowledge-use}部分的知识获取主题，这里的关键是每次新增标签具体减少了哪一种失配的不确定性\scite{ket-prabhu2021-clue}

跨观测空间还需要核对\emph{对应关系}。RGB图和深度图各自成像正常，不代表它们拍摄的是同一植株、同一时刻或同一视角。若用植物甲的RGB特征监督植物乙的深度图，两侧数据都可靠，监督关系却是错的。配对材料应检查对象身份、时间同步、视野覆盖、空间配准及缺测；桥接模态要能把两侧联系到同一对象或语义，共享类别名称也要确认粒度一致。只有教师行为时，能够对齐的是被查询输入上的响应，而非自动获得逐部位、逐像素的对应。

对新的负责人事实，证据核对同样分开进行：来源是否足以支持任命，发布日期与生效日期是否一致，机构身份是否匹配，任命是否只适用于某分支或代理期限。旧模型回答甲只能说明它曾经这样回答，不能否定当前任命。外部请求本身也不是无条件真值；未取得有效时间或主体身份时，可以保留待决并请求有针对性的补充证据。

撤回记录还需明确原始身份、重复副本、派生数据和已经影响的训练状态。若只知道文件名而不知道哪个数据快照用过它，删除该文件并不能确定重算范围；反之，不应把另一条独立合法记录仅因内容相似就自动列入撤回集合。没有充分证据时，补充监督、限制更新范围、暂缓自动更新和保留原版本都是可比较的选择。无标签输入无法识别任意条件规律变化，这一限制不能靠反复挑选“最符合旧模型”的样本消除，见第~\ref{thm:transfer-nonidentifiability}条定理。

\subsection{更新证据与验证证据的划分}
\label{sec:ket-update-split}

同一份材料如果既决定怎么更新，又用来宣布更新有效，就会掩盖对材料本身的适配。可将每条记录的用途写成“更新”“选择强度和停止时机”“最终验收”三种互斥标记，并记录分组身份、可见时间、标签到达时间和允许的批量处理方式。训练损失使用更新部分，约束系数与停止时机使用选择部分，最终验收部分不再用于反复挑选方案。记录的数量不同，不代表信息彼此独立。

例如，以$g_1,g_2,g_3$表示互不重叠的三个植株采集组，每组内的连拍整体分配用途，可以形成如下教学记录：
\[
\begin{array}{ccl}
g_1&\longrightarrow&\text{更新集：计算梯度}\\
g_2&\longrightarrow&\text{选择集：选约束强度与停止时机}\\
g_3&\longrightarrow&\text{验收集：候选冻结后评估}
\end{array}
\]
这里没有规定各集样本数，实际分配还要保证需要验证的环境和类别得到覆盖；关键是让用途跟随整组证据，而不是随单张文件任意改变。

同一植株的连拍应尽量按植株或采集事件分组，不能一张用于更新、近乎相同的另一张作为独立测试。同一事实的改写可以用于编辑训练，也可以用于泛化验证，但必须事先约定哪些模板、实体和事实族保留；已经用来选择编辑器的改写不再是独立的成功证据。连续时间窗口之间也可能共享环境状态或后续标签，应按照未来实际使用方式做时间切分，防止未来反馈倒流到过去的更新决策。

单次更新可以预先组织这些材料，连续更新则要在数据流中逐条维护用途。假设星期一的照片没有标签，系统先输出并进行获准的无标签适应；星期三标签才返回，那么星期一的即时预测必须以当时的状态计分，星期三的监督可以用于后续版本，不能重算星期一答案来替换当时结果。用于验证的延迟标签若后来转入训练，也应结束其独立验证身份，并另留后续验收证据。

测试时适应还要明确测试输入的使用权。允许一个批量共同估计归一化统计，与要求每个样本仅使用过去状态，是不同协议；允许反复遍历完整测试集，又与单遍流式适应不同。即使从未读取测试标签，联合处理未来测试输入也改变了信息条件。每个方案应记录批量大小、时间顺序、跨样本状态、重置边界和最终验证方式，使后文的更新计算和验证主张能够对应。

\section{知识的更新}
\label{sec:ket-update-operations}

证据和目标确定以后，需要选择改变系统的哪一部分。参数更新在固定结构内改变权重、偏置或校准量；结构调整改变组件、连接、路由与容量配置；外部更新改变可调用记录及其有效范围。三者可以组合：新增深度图编码器属于结构操作，学习其权重属于参数操作，为问答系统更新任命记录属于外部操作，而学习新的时间条件识别器又需要参数更新。

操作规模应服从变化范围。广泛的拍摄条件失配通常需要一批具有覆盖性的训练经验，单条高频时效事实可能更适合维护外部记录，指定参数编辑适用于可定位且能验证影响范围的请求。完整重训练可以重新估计大范围关系，但成本较高，也不会自动保留训练数据之外的有效编辑。局部微调、直接编辑或外部替换计算较便宜，同样不能据此推断已经满足关联一致性或训练影响消除要求。

\subsection{模型参数的更新}
\label{sec:ket-update-parameters}

固定模型结构以后，一次参数更新仍由多项选择共同决定：取用哪些经验、怎样赋权，构造什么学习目标、施加何种保留约束，允许改变哪些参数、从哪个版本开始，以及怎样执行和控制计算。下面沿这个构造过程展开。同一种回放材料可以进入监督损失或行为约束，同一个对齐目标也可以用于部分层或全模型更新；它们不是互斥的知识类别。定向编辑与共享基更新会跨越其中几步，下文在讨论更新范围时保留这两个完整案例，使位置选择、修改量和保留范围能够连在一起理解。

\subsubsection{训练经验的取用与加权}
\label{sec:ket-update-weighting}

原训练集中的频率不再代表当前环境时，平均计算每个样本的损失，实际优化的仍是原来的分布。若两域具有共同的输入输出含义，可以先问：能否不制造新样本，只改变已有样本对目标的贡献？\term{重要性加权}（importance weighting）把目标分布相对源分布更常见的情形赋予更高权重，使源样本上的平均损失能够估计目标风险。

设两域关于同一基准测度有联合密度$p_S,p_T$，损失$\ell(f_{\bm{\theta}}(\bm{x}),y)$可积，并满足目标被源覆盖，即$p_T(\bm{x},y)>0$的区域上$p_S(\bm{x},y)>0$。对离散标签，把目标风险中的每一项乘除源密度，可得
\begin{align}
R_T(\bm{\theta})
&=\sum_y\int \ell(f_{\bm{\theta}}(\bm{x}),y)
       p_T(\bm{x},y)\,\mathrm{d}\bm{x}\notag\\
&=\sum_y\int \ell(f_{\bm{\theta}}(\bm{x}),y)
       \frac{p_T(\bm{x},y)}{p_S(\bm{x},y)}
       p_S(\bm{x},y)\,\mathrm{d}\bm{x}\notag\\
&=\mathbb E_{(\bm{x},y)\sim\mathcal P_S}
  [w(\bm{x},y)\ell(f_{\bm{\theta}}(\bm{x}),y)].
\label{eq:ket-update-importance}
\end{align}
其中$w=p_T/p_S$是密度比。连续标签只需把求和改为积分。这个等式本身还不是可用算法，因为无标签目标样本一般不能确定联合密度比；偏移假设的作用，是把未知量缩减为可以估计的对象。

在协变量偏移下，条件标签概率相消，
\begin{equation}
\begin{aligned}
w(\bm{x},y)
&=\frac{p_T(\bm{x})p_T(y\mid\bm{x})}
       {p_S(\bm{x})p_S(y\mid\bm{x})}
=\frac{p_T(\bm{x})}{p_S(\bm{x})},\\
\hat R_w(\bm{\theta})
&=\frac{1}{n_S}\sum_{i=1}^{n_S}
\hat w(\bm{x}_{S,i})\ell(f_{\bm{\theta}}(\bm{x}_{S,i}),y_{S,i}).
\end{aligned}
\label{eq:ket-update-covariate}
\end{equation}
下标$S,i$表示第$i$个源样本。密度比可以直接估计，也可以利用区分来源和目标的域分类器间接估计：若训练该分类器时的域采样先验为$\pi_S,\pi_T$，其输出$d(\bm{x})$估计目标域概率，则贝叶斯公式给出
\[
\frac{p_T(\bm{x})}{p_S(\bm{x})}
=\frac{d(\bm{x})}{1-d(\bm{x})}\frac{\pi_S}{\pi_T}.
\]
这里的域概率必须相对于\emph{训练域分类器时}的混合分布定义；平衡采样与自然频率采样需要不同先验系数。权重应在参与最终验收之前估计，并检查估计误差，而不是把有限样本的域分类分数直接当作精确密度比。

作为教学计算，假设同一批花卉照片中晴天与阴影比例在源域是$0.8,0.2$，在目标域是$0.4,0.6$，且每一情形内部的输入分布和条件标签规律保持不变。一个固定识别器在两种情形下的平均损失分别为$0.1,0.3$。未校正的源平均是$0.8\times0.1+0.2\times0.3=0.14$；重要性权重为$0.4/0.8=0.5$和$0.6/0.2=3$，于是加权贡献分别为$0.8\times0.5\times0.1=0.04$、$0.2\times3\times0.3=0.18$，合计$0.22$，与直接按目标比例计算一致。这说明权重进入的是每条经验的损失项，不是随意把最终准确率乘一个比例。

大权重也会放大不确定性。令单个加权损失为$Z=w\ell$，独立源样本均值的方差为$\operatorname{Var}_S(Z)/n_S$；目标中常见但源中稀少的区域可能使这个方差很大。若损失非负且有界于$L$，截断为$w_c=\min(w,c)$后，
\[
0\le R_T-\mathbb E_S[w_c\ell]
=\mathbb E_S[(w-c)_+\ell]
\le L\,\mathbb E_S[(w-c)_+].
\]
这里$(u)_+=\max(u,0)$。截断限制了单样本影响，却引入偏差；再把截断权重归一化，还会改变上述偏差表达，不能继续宣称无偏。目标区域在源中完全缺失时，覆盖假设已经不成立，任何有限权重都不能补出从未观察过的标签信息。

标签偏移保留的是另一项条件关系。记源、目标类别先验分别为$\pi_{S,k}=p_S(y=k)$和$\pi_{T,k}=p_T(y=k)$，且目标出现的类别在源中都有非零概率，则
\begin{equation}
w(\bm{x},y=k)
=\frac{p_T(\bm{x}\mid k)\pi_{T,k}}
       {p_S(\bm{x}\mid k)\pi_{S,k}}
=\frac{\pi_{T,k}}{\pi_{S,k}}.
\label{eq:ket-update-label-weight}
\end{equation}
这时同一类别内部的所有样本使用同一权重。对两类花卉，假设源先验为$(0.8,0.2)$、目标先验为$(0.4,0.6)$，且类条件输入分布保持不变。若类内损失仍取教学数值$(0.1,0.3)$，校正后的两类贡献也是$0.04,0.18$。它与上一个算例具有相同算术，却依赖不同假设：这里按类别加权，前面按拍摄情形加权。不能只因数值相同就把两种偏移等同。

目标先验未知且没有目标标签时，可以利用固定预测器的混淆关系估计。令$\hat y=h(\bm{x})$，定义按真实类别归一化的混淆矩阵
\[
\bm C_{j,k}=\mathbb P_S(\hat y=j\mid y=k),
\qquad q_{T,j}=\mathbb P_T(\hat y=j).
\]
由于$h$固定，类条件输入分布不变推出$\mathbb P_T(\hat y=j\mid y=k)=\bm C_{j,k}$。全概率公式因而给出
\begin{equation}
\bm q_T=\bm C\bm\pi_T,
\qquad
\hat{\bm\pi}_T
=\argmin_{\bm\pi\ge0,\ \sum_k\pi_k=1}
\lVert\hat{\bm C}\bm\pi-\hat{\bm q}_T\rVert_2^2.
\label{eq:ket-update-bbse}
\end{equation}
源混淆矩阵需在与预测器训练适当隔离的有标签样本上估计，目标预测频率用允许访问的无标签样本估计。总体矩阵可逆时，方阵情形有$\bm\pi_T=\bm C^{-1}\bm q_T$；有限样本的直接求逆可能产生负比例或强烈噪声，带单纯形约束的估计至少保持概率合法。矩阵接近奇异时，微小频率误差仍会被放大。原始BBSE也可使用联合混淆矩阵估计先验比；这里使用列归一化形式，估计的是目标先验，两种写法的未知量不能混用\scite{ket-lipton2018-labelshift}

令教学混淆矩阵和目标预测频率为
\[
\bm C=
\begin{pmatrix}0.9&0.2\\0.1&0.8\end{pmatrix},
\qquad
\bm q_T=
\begin{pmatrix}0.48\\0.52\end{pmatrix}.
\]
由$0.48=0.9\pi_{T,1}+0.2(1-\pi_{T,1})$得到$\pi_{T,1}=0.4$、$\pi_{T,2}=0.6$，且$\det(\bm C)=0.7$非零。预测器并非完全准确，但不同真类产生不同的预测分布，足以恢复此例的先验。若两列完全相同，目标预测频率就不随类别比例改变，估计问题不可识别；把模型输出更自信地重复多次不能补上缺失信息。

校正还可以发生在预测端，而不重新训练底座。若源输出确实近似源后验，贝叶斯公式给出
\begin{equation}
p_T(k\mid\bm{x})
=\frac{p_S(k\mid\bm{x})\,\pi_{T,k}/\pi_{S,k}}
       {\sum_j p_S(j\mid\bm{x})\,\pi_{T,j}/\pi_{S,j}}.
\label{eq:ket-update-prior-correction}
\end{equation}
若源logit为$a_k(\bm{x})$，可以仅添加偏置$\log\pi_{T,k}-\log\pi_{S,k}$，编码器和原预测头权重保持不变。对源后验$(0.6,0.4)$，上例产生未归一化分数$(0.3,1.2)$，目标后验变为$(0.2,0.8)$。这里估计的是先验或校准系数，不是重新学会两类的特征。BBSE的先验识别不要求预测器概率校准，但式~\eqref{eq:ket-update-prior-correction}要将其当作源后验，因而依赖概率质量；这两项条件需要分开。类条件输入规律已变时，两种类别校正都没有一般保证。

经验权重也用于保留旧能力。设$S_{\mathrm{new}}$是当前更新数据，$M$是已经选定的历史记忆，\term{经验回放}（experience replay）把记忆中的样本再次送入训练，使旧情形继续产生梯度。一个明确的目标是
\begin{equation}
L_{\mathrm{mix}}(\bm{\theta})
=(1-\rho)\frac{1}{|S_{\mathrm{new}}|}
\sum_{(\bm{x},y)\in S_{\mathrm{new}}}\ell(f_{\bm{\theta}}(\bm{x}),y)
+\rho\frac{1}{|M|}\sum_{(\bm{x},y)\in M}
\ell(f_{\bm{\theta}}(\bm{x}),y),
\label{eq:ket-update-replay}
\end{equation}
其中$0\le\rho\le1$规定旧经验在优化中的份额，而不是直接等于旧样本数量比例。实际每步分别抽取新旧批量，计算两个批内均值后按系数相加；若直接拼接批量再平均，隐含权重就是批量大小比例。极小记忆上的回放研究表明，这种直接混合是需要认真比较的基线，其步骤包括取当前批量、抽取历史批量、共同更新以及更新记忆\scite{ket-chaudhry2019-replay}

在同一批100张新照片、容量20张的记忆下，假设新批量大小为10。每步加入10张回放照片、执行10步，总共处理100次新样本和100次旧样本；若每两个新批量之后才加入10张回放照片，共处理100次新样本和50次旧样本。两者记忆相同，计算量和旧经验权重却不同。若间隔回放改为每次20张，总旧样本处理次数可恢复到100，但梯度出现顺序、批量统计和中途漂移仍不相同。比较回放频率时应同时固定或报告这些量，不能只注明“使用同样大小的记忆”。

若总体希望优化离散目标$\sum_i a_i\ell_i$，其中$a_i\ge0$且$\sum_i a_i=1$，实际按$q_i>0$抽样，则单次无偏贡献是$a_i\ell_i/q_i$。不作这个校正，就在优化采样分布确定的另一个加权目标。刻意提高稀有花卉的作用完全可以合理，但应明确这是公平性、覆盖或风险取舍，而不是声称仍在估计原来的自然频率风险。无法被采到的记录也无法通过校正产生贡献。

历史标签并非唯一可回放的监督。DER保存样本和它在历史时刻的logit向量$\bm a_i^{\mathrm{old}}$，训练时使用
\begin{equation}
L_{\mathrm{DER}}(\bm{\theta})
=L_{\mathrm{new}}(\bm{\theta})
+\lambda_{\mathrm{logit}}\,
\mathbb E_{(\bm{x}_i,\bm a_i^{\mathrm{old}})\sim M}
\lVert\bm a_{\bm{\theta}}(\bm{x}_i)-\bm a_i^{\mathrm{old}}\rVert_2^2.
\label{eq:ket-update-der}
\end{equation}
缓存的预测来自经验流中不同历史时刻，不要求都来自任务结束的最优模型，因而需要记录对应的输出类别和版本。样本回放决定约束覆盖哪些输入，logit匹配决定希望在这些输入上保留什么行为；两者可以同时发生。DER++另外加入旧样本真实标签的监督损失，用于减轻错误历史预测被固定的问题，不能把该标签项无条件归入原始DER\scite{ket-buzzega2020-der}

生成回放用生成器产生旧输入，再配以旧标签或教师响应，形式上仍能进入式~\eqref{eq:ket-update-replay}或式~\eqref{eq:ket-update-der}，但训练目标反映的是生成分布。若生成器漏掉罕见阴影物种，反复生成常见照片无法保护该情形；若教师曾经答错，行为约束还可能保留错误。生成器本身继续训练时也会遗忘，因此要把生成质量、类别覆盖、教师版本和生成成本列入检查。

持续预训练中的经验单位通常是词元序列，连续指令学习中的单位则可能是指令与响应。新领域材料、历史材料和通用材料的混合比例决定模型用多少梯度继续巩固一般能力、用多少梯度吸收领域表达。按文档数等比例不等于按词元数等比例；长文档和重复指令可能占用更大训练份额。混合训练、先集中学习新域再回放旧域以及周期性插入通用数据，既改变训练顺序，也可能改变总计算和有效梯度，需同时检查领域收益、旧能力、新内容学习速度和总成本。

没有未来目标输入时，只能从已知条件中构造训练变化。\term{域随机化}（domain randomization）在仿真或可控生成过程中改变光照、纹理、视角等条件，使模型在训练时接触更多观测变化；Tobin等人的工作以仿真图像中的物体定位展示了这一路径。把它用于花卉图像时，随机化必须保持物种与标签的一致关系，不能把用于识别的花瓣结构随意改掉，也不能把渲染变化视为覆盖所有现实因素\scite{ket-tobin2017-randomization}

\term{领域泛化}（domain generalization）在不使用待部署目标域数据的协议下准备下一版本。若已知环境索引为$e$，稳健优化可以最小化$\max_e\hat R_e(\bm{\theta})$，避免训练平均值掩盖已知弱环境；但有限样本的最大经验风险可能被噪声支配，需要适当正则化，已知环境集合也不等于所有未来环境\scite{sagawa2020groupdro}

不变性方法则寻找在已知环境间保持有效的预测关系。例如，不变风险最小化要求共享表示上的同一个预测器在各训练环境都接近最优，而不只是让环境的输入统计相同。如果每个已知环境中背景颜色与物种都同样相关，这些数据就不能辨认背景是偶然关联，有限环境上的不变性仍可能选择错误特征。训练环境划分、假设和未来环境覆盖因此是必要条件；完整讨论见第~\ref{sec:domain-generalization}节\scite{arjovsky2019irm}

\subsubsection{更新目标与保留约束的构造}
\label{sec:ket-update-objective}

经验权重回答“哪条样本贡献多少”，却不能独自回答“应该调整哪一种关系”。雨季照片若仍包含辨认物种所需的形态，但编码器把阴影变化当成类别差异，就需要调整表示；旧任务若仍需服务，还要限制这种调整对历史预测的影响。设$f_{\bm{\theta}}=h_{\bm{\psi}}\circ\phi_{\bm{\omega}}$，其中$\phi$把输入变为表示，$h$据此预测，$\bm{\theta}=(\bm{\omega},\bm{\psi})$。下面逐项说明不同损失实际约束什么。

Deep CORAL通过匹配两域表示的协方差，减小特征各维之间的二阶统计差异。对域$D\in\{S,T\}$的一批$n_D>1$个表示，定义
\[
\bar{\bm z}_D=\frac1{n_D}\sum_i\bm z_{D,i},\qquad
\bm C_D=\frac1{n_D-1}\sum_i
(\bm z_{D,i}-\bar{\bm z}_D)(\bm z_{D,i}-\bar{\bm z}_D)^\top.
\]
若表示维数为$d$，其对齐损失为
\begin{equation}
L_{\mathrm{CORAL}}=\frac1{4d^2}\lVert\bm C_S-\bm C_T\rVert_{\mathrm F}^2.
\label{eq:ket-update-coral}
\end{equation}
训练时从两域取样，源标签形成任务损失，两域表示形成式~\eqref{eq:ket-update-coral}，两项共同更新编码器；使用时只需更新后的任务模型。损失匹配的是协方差，不是全部分布，也不单独匹配均值。若只有历史协方差缓存，还要保证其坐标与当前表示兼容；否则被比较的矩阵虽同型，却未必描述同一种特征\scite{ket-sun2016-coral}

DANN用另一个可训练模型辨认表示来自哪个域，再让编码器妨碍这种辨认。令域判别器$d_{\bm{\nu}}(\bm z)$输出目标域概率，域损失写为
\[
L_{\mathrm{dom}}
=-\mathbb E_S\log(1-d_{\bm{\nu}}(\phi_{\bm{\omega}}(\bm x)))
 -\mathbb E_T\log d_{\bm{\nu}}(\phi_{\bm{\omega}}(\bm x)).
\]
判别器更新$\bm{\nu}$以减小$L_{\mathrm{dom}}$，编码器更新$\bm{\omega}$以减小$L_{\mathrm{task}}-\lambda_{\mathrm{dom}}L_{\mathrm{dom}}$；预测头只接收相应任务信号。梯度反转层在前向时传递表示，在反向时把域损失对编码器的梯度乘以负系数，实现这两个相反方向。使用时域判别器通常不参与任务预测。与固定的协方差距离相比，域判别器能学习更复杂的分布差别，但“判别器分不清”还受其容量、训练程度和样本量影响\scite{ganin2016}

只对齐边际分布有一个直接的失败例子。假设两域各有两个等大的花卉簇，源域左簇是甲类、右簇是乙类，而目标编码器把甲类映到右边、乙类映到左边。两域表示分布完全一致，旧分类器却会把全部目标样本分错。把所有输入映到同一个常量也能消除域差别，却失去了分类信息。任务损失能够排除部分退化解，但不能仅凭源标签排除所有目标语义排列，相关不可识别性见定理~\ref{thm:transfer-nonidentifiability}。

条件对抗适应CDAN让域判别器同时看到表示$\bm z$与类别预测$\bm p$的关系。其多线性形式使用$\bm z\otimes\bm p$，即将每个特征分量与各类别概率相乘；判别器因而不仅比较“有多少这类特征”，还比较“这些特征被预测成什么类别”。训练仍由任务监督和域对抗构成，可按预测熵调整样本贡献，减少高不确定响应的影响。类别很多或表示很宽时，可用原文的随机多线性近似控制维数。若目标伪标签已经系统性反转，条件化可能强化错误对应；条件对齐提供新的约束形式，不自动提供缺失的真实类别\scite{ket-long2017-cdan}

需要显式建立样本间关系时，最优传输把对应写成质量分配。设源、目标样本的归一化权重为$\bm a,\bm b$，代价$\bm D_{i,j}$衡量将源表示$i$对应到目标表示$j$的代价，求
\begin{equation}
\min_{\bm\Gamma\ge0}
\sum_{i,j}\bm\Gamma_{i,j}\bm D_{i,j}
+\varepsilon\sum_{i,j}\bm\Gamma_{i,j}\log\bm\Gamma_{i,j},
\quad
\bm\Gamma\bm 1=\bm a,\quad
\bm\Gamma^\top\bm1=\bm b.
\label{eq:ket-update-transport}
\end{equation}
这里$\bm\Gamma$是运输耦合，允许一个源样本向多个目标样本分配质量，$\varepsilon\ge0$调节熵正则，约定$0\log0=0$。平方距离代价下，源点的重心映射为$\tilde{\bm z}_{S,i}=\sum_j\bm\Gamma_{i,j}\bm z_{T,j}/a_i$；可在映射后的带标签源样本上估计预测器，也可把传输代价接入可学习表示。类别或邻域正则还能限制不合理的混合。测试新输入时必须说明使用原目标坐标上的预测器，还是另学了可外推映射，有限样本的耦合矩阵本身不是任意新输入的编码器。代价中的几何接近若不反映物种语义，最优耦合也可能是错误对应\scite{ket-courty2015-transport}

这些目标都需要区分“去掉无关差异”和“抹掉有用信息”。例如，花瓣颜色可能既受相机影响又有物种判别价值，不能仅因它能预测域身份就将其全部删除。对齐强度应由允许的验证证据选择；有当前标签时，要直接检查类内和类间关系，没有标签时则应承认统计对齐只能支持有限的代理判断。

观测空间改变后，对齐对象还要重新定义。跨模态蒸馏可以让冻结的RGB教师为同一对象的深度图提供中间表示目标。设配对数据为$(\bm x_{A,i},\bm x_{B,i})$，$A$是旧模态、$B$是新模态，可学习
\begin{equation}
L_{\mathrm{pair}}
=\frac1n\sum_i
\left\lVert g_{\bm\eta}(\phi_{B,\bm\omega}(\bm x_{B,i}))
-\phi_{A,\mathrm{old}}(\bm x_{A,i})\right\rVert_2^2.
\label{eq:ket-update-crossmodal}
\end{equation}
$g_{\bm\eta}$把新表示转换到可比较的空间，不要求两种原输入维数相同。训练时教师和配对旧观测可见，部署时只使用新观测、学生及所需映射。原始跨模态蒸馏以配对RGB与深度或光流传递监督，关键依据是对象对应，而不是两批数据都来自同一个大类别\scite{ket-gupta2015-crossmodal}

广义蒸馏把训练阶段独有的信息称为特权信息：教师可以看更丰富的输入，学生只能看部署时拥有的输入，学生同时学习真实标签与教师软预测。这样的教师可能给出比硬标签更细的学习信号，但不能消除学生输入本身的不确定性。若两种仅以花色区分的物种在深度图中完全同分布，任意确定映射都不能为每张深度图恢复真实花色；平方损失下的最佳映射至多趋向给定深度后的教师表示条件均值。共同语义空间也必须由配对、共享标签或可靠桥接关系建立，不能仅凭两套向量都具有$d$维便宣布语义统一。蒸馏的一般训练过程见第~\ref{sec:ket-integration-student-training}节，多模态信息的取得属于第~\ref{part:efficient-knowledge-use}部分的内容范围\scite{ket-lopezpaz2016generalized}

跨语言更新可作同样检查。平行表达和实体身份能够把不同词序、词形的句子连接到共同含义，但专名翻译、日期表达和否定范围仍需核对。若同一个多语言编码器已经形成可用的共享语义，主要问题只是某种语言的材料变少，可以先调整经验混合和权重；若新语言缺少关键表达的覆盖，则需要补充对应监督。增加新语言的输出规范或改变成功标准时，任务要求交第~\ref{sec:ket-transfer-output-semantics}节确定，不能只按输入分布变化处理。

换成回归、时序或图数据以后，前述关系仍需逐项验证。回归对齐要保留数值单位、量纲与连续目标的条件关系，离散类别先验校正不能原样套用；将厘米误当米不是通过表示分布匹配就能合理修复的问题。时序样本具有相邻依赖，更新时只能使用当时可得的历史窗口，不能用未来标签选择对齐方向；打乱批量还可能改变任务需要的时间关系。图模型中的一个节点表示取决于邻居、边类型及消息传播，更新节点特征或图结构会影响多个节点，跨图对应需说明实体和关系含义。分类上的边际对齐结果不为这些新增条件提供自动保证。

当当前证据已经支持新规律时，保留机制要防止对其余有效能力的无谓改动。最直接的参数约束是让$\bm\theta$靠近固定历史参考$\bm\theta^\star$；但所有参数同等受罚，忽略了不同方向对旧任务的影响。\term{弹性权重巩固}（elastic weight consolidation, EWC）按旧任务的局部敏感程度赋权，使重要方向较难移动。下面从历史后验说明这项二次惩罚从何而来。

设$S_{\mathrm{old}},S_{\mathrm{new}}$在给定参数后满足相应的似然分解，新阶段的负对数后验去掉常数后为
\[
-\log p(\bm\theta\mid S_{\mathrm{old}},S_{\mathrm{new}})
=-\log p(S_{\mathrm{new}}\mid\bm\theta)
-\log p(\bm\theta\mid S_{\mathrm{old}})+\mathrm{const}.
\]
在历史后验的局部极大点$\bm\theta^\star$附近，一阶项为零。记历史负对数后验的Hessian为$\bm H_{\mathrm{old}}$，局部二阶展开给出
\[
-\log p(\bm\theta\mid S_{\mathrm{old}})
\approx\mathrm{const}
+\tfrac12(\bm\theta-\bm\theta^\star)^\top
\bm H_{\mathrm{old}}(\bm\theta-\bm\theta^\star).
\]
在适当正则条件下，模型分布下负对数似然的期望曲率可用Fisher信息表示。EWC用其对角近似构造可存储的权重：
\begin{align}
F_j&=\mathbb E_{\bm x\sim\mathcal P_{\mathrm{old}},\
y\sim p_{\bm\theta^\star}(\cdot\mid\bm x)}
\left[
\left(\frac{\partial}{\partial\theta_j}
\log p_{\bm\theta^\star}(y\mid\bm x)\right)^2
\right],\notag\\
L_{\mathrm{EWC}}(\bm\theta)
&=L_{\mathrm{new}}(\bm\theta)
+\frac{\lambda_{\mathrm E}}2
\sum_j\hat F_j(\theta_j-\theta_j^\star)^2.
\label{eq:ket-update-ewc}
\end{align}
这里$\mathcal P_{\mathrm{old}}$是历史输入的分布，标签按参考模型的条件分布抽样。样本数、平均损失尺度和先验近似的系数在实践中常并入$\lambda_{\mathrm E}$。若用观测标签代替模型采样标签，得到的是常称的经验Fisher估计，并不无条件等于上式Fisher或有限样本Hessian。对角化还舍弃了参数间的相关方向，所以式~\eqref{eq:ket-update-ewc}是局部近似的保留机制，不是完整历史后验\scite{kirkpatrick2017ewc}

参考参数和$\hat F_j$应在新阶段开始前根据历史状态固定。以后不保存旧原样本仍可计算二次罚项，但先前已经使用历史数据或等价依据估计了重要性，并保存两份参数规模的向量。对参数$j$，保留项的梯度是$\lambda_{\mathrm E}\hat F_j(\theta_j-\theta_j^\star)$，它把离开参考的参数拉回；若参考也随每一步一起移动，约束就不再保护同一个历史状态。原始任务式EWC在任务结束处确定参照，多个任务的近似如何累积或合并也需声明，不能在没有边界的流上默认获得任务分段。

SI使用训练轨迹而不是在终点重新测量局部敏感性。对一次已经结束的任务学习区间，记任务损失梯度为$g_j^{(s)}$，实际参数位移为$\Delta\theta_j^{(s)}$。下面采用显式截去负估计的非负化写法：
\[
\omega_j=-\sum_s g_j^{(s)}\Delta\theta_j^{(s)},\qquad
\Omega_j\leftarrow\Omega_j+
\frac{\max(\omega_j,0)}{(\theta_{j,\mathrm{end}}-\theta_{j,\mathrm{start}})^2+\xi},
\]
其中$\xi>0$避免小净位移造成除零，$\Omega_j$初始为零。$\omega_j$是离散路径积分近似：某参数的移动若持续降低任务损失，就得到较大的贡献；随机梯度、优化器和路径都会影响估计。将负估计截为零是这里声明的数值处理，不把有限步轨迹无条件当成非负重要性。任务结束后固定$\Omega_j$和新的参考，再用加权平方偏离限制下一阶段。SI不需要为了这项累积额外保存全部旧样本，但需要追踪训练路径和确定巩固时机；它与EWC衡量的是不同近似下的重要性，都不能将某条事实唯一定位到一个参数\scite{ket-zenke2017-si}

若希望限制的是旧回答而不是参数距离，可以直接匹配历史行为。令$U$为允许查询且需要保留的输入集合，冻结旧模型后，以温度$T>0$计算
\[
q_{i,k}=\frac{\exp(a^{\mathrm{old}}_{i,k}/T)}
{\sum_j\exp(a^{\mathrm{old}}_{i,j}/T)},\qquad
p_{i,k}(\bm\theta)=\frac{\exp(a_{i,k}(\bm\theta)/T)}
{\sum_j\exp(a_{i,j}(\bm\theta)/T)}.
\]
类别顺序和被保留的输出范围在两侧保持一致。将$U$上各输入、各类别的固定旧概率合记为$\bm q$，以它们为参照的行为损失是
\begin{align}
L_{\mathrm{beh}}(\bm\theta)
&=\frac{T^2}{|U|}\sum_{i\in U}\sum_k
q_{i,k}\log\frac{q_{i,k}}{p_{i,k}(\bm\theta)}\notag\\
&=\mathrm{const}
-\frac{T^2}{|U|}\sum_{i\in U}\sum_k q_{i,k}\log p_{i,k}(\bm\theta).
\label{eq:ket-update-behavior}
\end{align}
第二行利用了$q$固定；它与软标签交叉熵仅差一个不随学生变化的常数。对批平均损失$L_{\mathrm{beh}}$中的单个logit$a_{i,k}$求导，得到$\frac{T}{|U|}(p_{i,k}-q_{i,k})$，说明更新在缩小概率分布的差别。只对一个样本的logit求导，仍须保留外层平均因子。前面的$T^2$抵消了一部分温度造成的梯度尺度变化，但不能让不同温度下的训练完全等价。

若保存了完整softmax概率并知道原温度$T_0$，即使没有原始logit，也可以将各概率取$T_0/T_1$次幂后重新归一化，得到新温度$T_1>0$下的分布；logit共同的加性常数会在归一化中消去。仅保存硬类别、截断或低精度概率，或者不知道原温度及接口后处理方式时，则不能保证精确重建原模型在指定温度下的预测。

LwF在当前任务输入上查询冻结旧模型的旧任务输出，以该行为损失配合当前监督学习。它不要求旧任务原样本回放，却仍需要旧模型或提前缓存的旧预测；当前输入若远离旧任务的有效区域，就不能充分约束那里未被查询的行为。使用时保留更新后的网络和所需任务头，不必让教师一直在线。与DER相比，LwF的典型约束输入是当前数据，DER则明确保存并回放历史输入及轨迹预测，这一区别直接改变覆盖范围\scite{ket-li2016-lwf}

在没有即时标签时，保留对象也可以是表示。CaSSLe冻结历史编码器，让当前编码器的表示通过一个可学习预测器去匹配旧表示，同时继续优化当前自监督目标；匹配沿用对应自监督方法的相似性或对比损失，不要求所有方法都使用简单平方距离。预测器允许当前表示作有用变化，而仍能读出旧表示所包含的信息。旧编码器与当前输入提供保留信号，标签可以稀疏或以后才到来，但缺少旧输入覆盖的问题仍存在。事后收集标签训练线性探测器，回答的是表示中能否读出某种信息；没有相应预测头时，不能把该分数当作系统当时已经会作出的任务回答\scite{ket-fini2021-cassle}

把这些选择放入同一次更新，可以写出一个供分析的共同目标
\begin{equation}
\begin{split}
L(\bm\theta)={}&L_{\mathrm{new}}(\bm\theta)
+\lambda_{\mathrm R}L_{\mathrm{replay}}(\bm\theta)
+\lambda_{\mathrm A}L_{\mathrm{align}}(\bm\theta)\\
&+\lambda_{\mathrm B}L_{\mathrm{beh}}(\bm\theta)
+\frac{\lambda_{\mathrm P}}2
\sum_j\Omega_j(\theta_j-\theta_j^\star)^2.
\end{split}
\label{eq:ket-update-joint-objective}
\end{equation}
这里$L_{\mathrm{replay}}$是历史记忆上的任务损失，$\lambda_{\mathrm R}\geq0$表示相对于系数为1的新损失的回放权重，与式~\eqref{eq:ket-update-replay}中的归一化份额$\rho$不同。$L_{\mathrm{align}}$由当前可用的两域或配对证据构造，$L_{\mathrm{beh}}$只作用于有效保留范围，$\Omega_j$可取已说明的重要性估计。不是所有应用都应启用全部项；域对抗还需要另行更新判别器，不能将其两方优化误写成所有参数共同下降同一个标量。即使用同一组新旧样本，统计对齐约束的是表示关系，行为匹配约束的是输出，二次罚项约束的是参数位移，三者并不等价。

目标之间可能存在真实冲突。对于2026年7月15日的甲所负责人，修改损失要求乙，若行为损失同时要求严格保持旧答案甲，就不存在同时满足两项硬要求的预测器。增大两个损失的系数不会解决语义冲突，只会改变折中结果。应先从保留集合中移除已被合法修正的旧要求，保留6月历史查询和无关实体；在允许软折中的任务上，再用验证数据选择约束强度。图~\ref{fig:ket-update-objective}标出了固定参照和随更新变化的量，便于检查是否误把当前模型的输出作为自己必须保持的历史证据。

\begin{figure*}[htbp]
\centering
% All losses refer to the same current parameter vector.
% Fixed teacher and parameter references must not move with the student.
\begin{tikzpicture}[
  x=1mm,y=1mm,font=\figurefont\small,
  box/.style={rectangle,draw=BookInk,fill=BookPaper,line width=0.7pt,
    align=center,inner xsep=2mm,inner ysep=2mm,minimum height=12mm},
  flow/.style={draw=BookInk,line width=0.8pt,-{Latex[length=2mm]},
    rounded corners=1mm},
  constraint/.style={flow,dashed,draw=BookBlue},
  source/.style={box,text width=37mm},
  computed/.style={box,text width=32mm,fill=FigureOutput!10},
  loss/.style={box,fill=FigureLoss!10,text width=24mm}
]
  \node[anchor=west,text=BookInk] at (1,98)
    {各行共用当前模型$f_{\bm\theta}$，历史参照在本轮固定};
  \node[source,fill=FigureInput!10] (data) at (22,80)
    {当前与回放批量\\输入及有效标签};
  \node[source,fill=FigureInput!10] (pairs) at (22,56)
    {两域或配对输入\\明确对应依据};
  \node[source,fill=FigureModel!10] (teacher) at (22,32)
    {固定旧模型与输入$U$\\提供历史预测$\bm q$};
  \node[source,fill=FigureModel!10] (reference) at (22,8)
    {固定参考$\bm\theta^\star$\\固定重要性$\bm\Omega$};
  \node[computed] (prediction) at (70,80)
    {当前预测\\$f_{\bm\theta}(\bm x)$};
  \node[computed] (features) at (70,56)
    {当前表示\\$\phi_{\bm\omega}(\bm x)$};
  \node[computed] (behavior-pair) at (70,32)
    {当前$\bm p_{\bm\theta}(U)$\\与固定$\bm q$比较};
  \node[computed] (difference) at (70,8)
    {当前参数偏离\\$\bm\theta-\bm\theta^\star$};
  \node[loss] (task) at (112,80) {任务损失\\监督与回放};
  \node[loss] (align) at (112,56) {对齐损失\\统计或对应};
  \node[loss] (behavior) at (112,32) {行为损失\\概率匹配};
  \node[loss] (penalty) at (112,8) {参数惩罚\\重要性加权};
  \node[box,text width=24mm,fill=FigureOutput!10] (update) at (151,44)
    {组合有效目标\\同一次\\参数更新};
  \draw[flow] (data) -- (prediction);
  \draw[flow] (pairs) -- (features);
  \draw[constraint] (teacher) -- (behavior-pair);
  \draw[constraint] (reference) -- (difference);
  \draw[flow] (prediction) -- (task);
  \draw[flow] (features) -- (align);
  \draw[flow] (behavior-pair) -- (behavior);
  \draw[flow] (difference) -- (penalty);
  \draw[BookInk,line width=0.8pt] (132,8) -- (132,80);
  \draw[BookInk,line width=0.8pt] (task.east) -- (132,80);
  \draw[BookInk,line width=0.8pt] (align.east) -- (132,56);
  \draw[BookInk,line width=0.8pt] (behavior.east) -- (132,32);
  \draw[BookInk,line width=0.8pt] (penalty.east) -- (132,8);
  \draw[flow] (132,44) -- (update.west);
\end{tikzpicture}

\caption{新旧证据形成同一次参数更新的目标。实线表示输入、表示或预测的计算流，虚线表示固定参照对损失的约束。历史记忆可同时提供回放标签和旧行为，因此回放与蒸馏并不互斥；参考参数及重要性在本轮更新前固定，当前模型的表示、预测和参数随优化变化。图中并列损失是可选组合，不要求每次更新全部启用。}
\label{fig:ket-update-objective}
\end{figure*}

\FloatBarrier
\subsubsection{更新范围与起点的确定}
\label{sec:ket-update-range}

学习目标相同，允许变化的参数不同，能够到达的解也不同。若当前表示仍能区分两类花卉，只是类别先验变了，可以只估计式~\eqref{eq:ket-update-prior-correction}的偏置；若边界方向变了，可以重新学习预测头；若编码器已经丢弃新环境中必要的纹理，则只改预测头不能恢复这些信息。选择部分层、全部参数或附加的少量可训练参数，需要同时说明表达能力、旧能力风险和计算预算。参数范围的取舍见第~\ref{sec:ket-transfer-parameter-adaptation}节，附加参数的计算位置与机制见第~\ref{sec:ket-transfer-added-parameters}节；这里着重说明哪些参数允许改变、从哪个版本开始，以及这两项选择怎样影响保留范围。

起点同样是方法的一部分。现有版本通常保留最近学到的内容，较早快照可能避开已经发生的失稳，重新初始化则放弃旧参数提供的起点。三者需要重估的内容、训练代价和有效请求继承情况都不同。允许更新的参数集合$\mathcal J$确定以后，可以写$\theta_j=\theta_j^\star$对所有$j\notin\mathcal J$成立；这只限制可达参数空间，并未声明其中每个可达模型都能保持旧能力。恢复快照还需要核对该版本之后已经生效的编辑与撤回，不能先上线再补查。

定向事实编辑把更新范围缩到较少请求，但仍包含三个职责：根据请求选择候选位置，计算参数修改量，以及验证实际语义影响。局部微调直接在选定层上最小化编辑损失，并以保留样本或位移约束限制外溢；其位置选择可以来自先验经验或诊断，不能因损失下降就反过来证明定位正确。学习编辑器则预先训练一个“从请求到参数增量”的模型，将反复求解的部分计算摊到编辑器训练中。

MEND利用全连接层梯度的结构。对一条输入，权重梯度常可写成误差信号与层输入的外积$\bm u\bm v^\top$；长序列或批量对应这些外积之和。编辑器分别处理这些低秩因子，再组合成所选层的修改量，避免为完整权重矩阵构造同样巨大的编辑网络。训练编辑器时，从编辑请求构造内层修改，再用修改后模型在目标请求和局部性样本上的损失更新编辑器；使用时固定编辑器，为新请求计算梯度因子和一次参数增量。它需要具有代表性的编辑训练分布，编辑器训练、所需基座访问与每次请求成本应分开核算，不能只统计应用增量的时间\scite{ket-mitchell2021-mend}

ROME采用定位后写入的路径：通过干预与恢复隐藏状态的分析寻找与事实回忆有关的位置，再修改选定MLP中的线性映射。把待改矩阵记为$\bm W$，其输入激活$\bm k$称为键，输出激活$\bm v=\bm W\bm k$称为值。针对编辑主体取得的键$\bm k_\star$依赖主体、上下文和选定层；值是网络内部向量，不是答案文本。为写入新答案，先把该位置的输出替换为可优化的向量，让模型继续完成后续计算，根据新答案的预测损失及保留要求求得目标激活向量$\bm v_\star$，再修改$\bm W$使指定键映射到它。在$\bm C=\mathbb E[\bm k\bm k^\top]$正定的简化条件下，保留其他键响应的加权最小修改可写为
\begin{equation}
\min_{\Delta\bm W}\operatorname{tr}
(\Delta\bm W\bm C\Delta\bm W^\top),
\qquad
(\bm W+\Delta\bm W)\bm k_\star=\bm v_\star.
\label{eq:ket-update-rome-problem}
\end{equation}
迹罚项等于$\mathbb E\|\Delta\bm W\bm k\|_2^2$，衡量旧键分布上线性响应的平均平方变化。令残差$\bm r=\bm v_\star-\bm W\bm k_\star$。逐行求解带线性约束的二次最小化，得到
\begin{equation}
\Delta\bm W=
\frac{\bm r(\bm C^{-1}\bm k_\star)^\top}
{\bm k_\star^\top\bm C^{-1}\bm k_\star}.
\label{eq:ket-update-rome}
\end{equation}
代入即可验证$\Delta\bm W\bm k_\star=\bm r$，且修改秩至多为一。$\bm C$是未中心化二阶统计，不应与CORAL的中心化协方差混用；奇异或病态时需要明确正则化等数值处理。该局部线性关系仍嵌在非线性网络内，满足写入约束不能代替改写、历史语义和多跳关系验证\scite{meng2022rome}

MEMIT面向批量事实写入，将多个请求的键组成$\bm K$、待补偿残差组成$\bm R$，并在多个关键MLP层分配残差。在一个层级的正则化最小二乘写法中，
\[
\min_{\Delta\bm W}
\lVert\Delta\bm W\bm K-\bm R\rVert_{\mathrm F}^2
+\operatorname{tr}(\Delta\bm W\bm C\Delta\bm W^\top)
\]
给出$\Delta\bm W=\bm R\bm K^\top
(\bm C+\bm K\bm K^\top)^{-1}$；这里$\bm C$已包含旧键统计的保留权重。与单条硬写入不同，这个式子在批量拟合与旧响应保持之间折中，相关或冲突的键会相互影响。算法逐层写入后重新计算下游键，并重新分配尚未实现的残差，不能把同一增量不加区分地复制到每一层。请求批量越大，键、残差、统计估计与线性求解的资源负担越重要；一次批量编辑与逐条连续编辑也不是相同的训练轨迹\scite{ket-meng2022-memit}

限制可训练参数、限制改动范数和限制无关响应，是三个不同层次。若共享标量参数同时影响温室和雨季两个预测头，改动这一个参数就可能改变两个任务；若大量参数作相互抵消的重参数化，参数变化很多，函数却可能不变。因此，参数数量不能直接衡量语义局部性。位移小也不保证接近决策边界的样本不换类，真正关心的响应限制要在相应输入上建立并验证。下一小节的二维算例将直接计算这一区别。

持续预训练和任务微调也需要这项范围判断。只调整输出层可能保留一般表示，却限制领域知识的吸收；全模型训练允许更大变化，也需要通用任务和历史语言能力的证据约束。连续更新时应记录每轮可训练层、学习率及损失比例，避免把训练词元增加、领域语料质量改善和范围扩大造成的收益混为一谈。旧能力可能由多个层共同实现，不能假定底层冻结便保护了所有通用能力。

类别扩展还会改变共同预测器的竞争关系。BiC在主学习阶段用新数据、旧样本及蒸馏学习分类器，随后固定已有网络，在单独保留、旧新类别适当平衡的验证材料上学习新类logit的仿射校正$\tilde a_k=\alpha a_k+\beta$，旧类logit不作同样变换。第二阶段的目的，是缓解新旧训练份额不均形成的预测偏差；它依赖代表性的校正材料，也消耗本可用于主训练的样本。严格保持旧类之间的相对响应，不保证加入新类后的整体决策不变，校准与保留需要共同检查。扩展哪些类别由第~\ref{sec:ket-transfer-output-semantics}节确定，这里处理扩展后的能力竞争\scite{ket-wu2019-bic}

即使当前任务和旧任务暂时都能完成，系统也可能越来越难吸收后续经验。长期训练可能出现长期不活跃的单元、过大的权重或优化状态不适配等现象，是否造成可塑性下降要以同预算的新内容学习检验，而不能只看一种内部统计。持续反向传播在原有连接不变时，跟踪单元效用和成熟程度，选择部分低效且足够成熟的单元，将其输入权重重新初始化、输出权重置零，并重置相关统计；输出置零限制了新随机特征的即时影响，成熟期避免新单元刚生成就被再次替换。这是权重重置而非网络扩容，其原始研究也明确不以解决遗忘为保证：当前低效单元可能承载稀有旧能力，重置前后仍要复测\scite{ket-dohare2024-plasticity}

可塑性还与共享参数如何吸收经验有关。ELLA把任务参数写作$\bm\theta_k=\bm L\bm s_k$，共享基$\bm L\in\mathbb R^{d\times r}$提供可组合方向，稀疏系数$\bm s_k$说明任务使用哪些方向。在基维数$r$和连接固定时，新任务到来先取得任务自身的参数估计$\hat{\bm\theta}_k$及局部二次损失权重$\bm D_k$。把损失归一化系数计入$\bm D_k$，可通过
\[
\min_{\bm s}
(\hat{\bm\theta}_k-\bm L\bm s)^\top
\bm D_k(\hat{\bm\theta}_k-\bm L\bm s)+\mu\lVert\bm s\rVert_1
\]
估计当前任务的稀疏系数，再固定各任务最近取得的系数，利用二次目标的累积统计更新共享基，并加上基的正则化。新任务改变$\bm L$以后，旧任务即使不重新估计系数，其预测参数$\bm L\bm s_k$也会变化。局部二次近似和固定旧系数减少了反复访问全部旧样本与重新求解全部任务的成本，代价是需要保存相应参数、曲率信息或累积统计；这些近似不等于原联合目标被精确求解\scite{ket-ruvolo2013-ella}

在这个过程中发生的是参数共享和精炼，而不是自动增加结构容量。使用时需取得相应任务系数，未知任务身份不能凭空选中正确系数。历史经验若让后续任务以更少样本得到可用系数，才支持前向收益；基的改变若改善旧任务，才支持后向收益，两者都需对照而非由分解形式直接推出。

\subsubsection{更新计算的实施与控制}
\label{sec:ket-update-computation}

目标、证据和范围已经明确，才可以把它们接成一次可执行更新。阶段开始时保存仍需固定的教师、参数参照、重要性和请求范围；每步按约定组成新旧批量，计算当前模型的任务输出与表示，再计算候选梯度或直接修改量。优化器只改变允许更新的参数，必要时投影方向或拒绝候选步。生成候选版本后，还要用未参与这一选择的证据验收。优化目标下降只证明这个目标在所用材料上下降，不能跳过语义和任务验证。

普通梯度步骤是$\bm\theta^{(t+1)}=\bm\theta^{(t)}-\eta\bm g^{(t)}$。受约束步骤先把$\bm g^{(t)}$替换为满足保留条件的方向，直接编辑则可能利用式~\eqref{eq:ket-update-rome}一次计算参数增量。三者都需要记录实际改动层和数值，计算成本却不同：普通训练重复前后向传播，约束方法另算旧任务梯度并求解子问题，局部编辑还可能需要定位、统计估计或预先训练编辑器。保存动量、二阶矩和随机状态能帮助复现训练，但这些状态也占资源，并可能包含撤回请求涉及的历史信息。

考虑指定的二维教学算例。线性模型$f_{\bm w}(\bm x)=\bm w^\top\bm x$原权重为$\bm w=(1,1)^\top$，请求把$\bm x_a=(1,0)^\top$的输出从1改为2。令增量为$\Delta\bm w$，最小改动问题为
\[
\min_{\Delta\bm w}\lVert\Delta\bm w\rVert_2^2,
\qquad
(\bm w+\Delta\bm w)^\top\bm x_a=2.
\]
约束给出$\Delta w_1=1$，未受约束的第二分量取0使范数最小，所以$\Delta\bm w=(1,0)^\top$，新权重为$(2,1)^\top$。输入$\bm x_b=(0,1)^\top$的输出仍为1，但$\bm x_c=(1,1)^\top$的输出由2变成3。只改一个权重保住了一个非目标响应，却改变了另一个非目标响应。

若另有明确要求保护$\bm x_c$的输出2，就增加$\Delta w_1+\Delta w_2=0$。与$\Delta w_1=1$联立，得到$\Delta\bm w=(1,-1)^\top$和新权重$(2,0)^\top$，平方改动代价由1增至2，欧氏范数由1增至$\sqrt2$。新约束保住了$\bm x_c$，却使$\bm x_b$输出由1变为0；若三个响应都要按上述要求满足，此二维模型便没有足够自由度。更直接地，对同一个$\bm x_a$同时要求输出2和保持1，不论扩大优化步数还是调整范数，都没有可行解。该算例解释约束的几何关系，不替代ROME等方法的层定位、键值估计和二阶统计处理。

对多个旧任务，不一定能写出简单的线性输出约束，可以改用历史记忆上的损失变化。设旧任务$k$的记忆损失为$L_k(\bm\theta)$，其当前梯度$\bm g_k=\nabla L_k(\bm\theta)$，候选下降方向为$\bm v$，实际步长为$-\eta\bm v$。Taylor展开给出
\begin{equation}
L_k(\bm\theta-\eta\bm v)
=L_k(\bm\theta)-\eta\bm g_k^\top\bm v
+O(\eta^2\lVert\bm v\rVert_2^2).
\label{eq:ket-update-gradient-taylor}
\end{equation}
因此，一阶近似下不增大旧损失的条件是$\bm g_k^\top\bm v\ge0$，不是$\le0$；梯度与下降位移的负号决定了方向。GEM从当前任务梯度$\bm g$出发，求最接近它的可行方向：
\begin{equation}
\min_{\bm v}\frac12\lVert\bm v-\bm g\rVert_2^2,
\qquad \bm G^\top\bm v\ge\bm0,
\quad
\bm G=(\bm g_1,\ldots,\bm g_K).
\label{eq:ket-update-gem}
\end{equation}
若原方向已满足全部条件，直接使用$\bm g$；否则求解这个凸二次规划。每个旧任务的记忆只提供相应经验损失梯度，并非总体风险的精确梯度\scite{ket-lopezpaz2017-gem}

参数维数可能远大于旧任务数，可以转到任务维数的对偶。对不等式引入$\bm\alpha\ge0$，Lagrange函数为
\[
\mathcal L(\bm v,\bm\alpha)
=\tfrac12\lVert\bm v-\bm g\rVert_2^2
-\bm\alpha^\top\bm G^\top\bm v.
\]
对$\bm v$求驻点得到$\bm v=\bm g+\bm G\bm\alpha$，代回后等价于
\begin{equation}
\min_{\bm\alpha\ge0}
\tfrac12\bm\alpha^\top\bm G^\top\bm G\bm\alpha
+\bm\alpha^\top\bm G^\top\bm g,
\qquad
\bm v^\star=\bm g+\bm G\bm\alpha^\star.
\label{eq:ket-update-gem-dual}
\end{equation}
例如$\bm g=(-1,1)^\top$、$\bm g_1=(1,0)^\top$，原内积为$-1$。约束要求$v_1\ge0$，投影得到$\bm v^\star=(0,1)^\top$，对应$\alpha^\star=1$；当前方向中会一阶损害旧任务的分量被去掉，但当前任务的下降速度也可能减小。

一阶可行不意味着有限步长安全。若$L_k$的梯度是$\beta_k$-Lipschitz，则
\[
L_k(\bm\theta-\eta\bm v)
\le L_k(\bm\theta)-\eta\bm g_k^\top\bm v
+\frac{\beta_k\eta^2}{2}\lVert\bm v\rVert_2^2.
\]
要由这个上界推出不增，需内积至少为$\beta_k\eta\lVert\bm v\rVert_2^2/2$；内积为零时，二阶项仍可能使损失上升。随机记忆梯度、求解误差和曲率变化又增加偏差，所以可减小步长、复算记忆损失或退回候选步，但不能宣称整段训练轨迹都无遗忘。若后续使用动量或自适应预条件改变了实际位移，单独投影原始梯度也未必约束住实际步骤，约束对象须与优化器输出一致。

MER把迁移与干扰联系到梯度相容性。对经验$i,j$，先沿$\bm g_i$走一小步，会使另一个损失近似变为$L_j(\bm\theta)-\eta\bm g_j^\top\bm g_i$；正内积有利于相互降低损失，负内积提示干扰。MER从当前样本与回放记忆构造小批量，执行内部学习，再在批内、批间将参数向学习后的方向作Reptile式插值，通过这种元学习过程近似鼓励跨经验梯度相容。它不是每步检查全部旧任务并投影的GEM硬约束，内积也是当前参数附近的局部量。减少负内积可能改善共同学习，但不保证将来未见任务更容易；前向、后向收益仍需分别测量\scite{ket-riemer2018-mer}

约束太强时，可行方向可能几乎为零，使新内容无法学入。允许放松的应是已确认过时或可以折中的要求，而不是事后删去所有受损任务来改善分数；事实修正和撤回更要按有效范围选择约束。优化状态重置有时能改变随后几步的方向，但不直接恢复参数中缺失的信息。放松强度、改变步长、扩大参数范围和增加证据针对的是不同限制，选择哪一种应由失败诊断决定。

使用阶段也可能允许更新。\term{测试时适应}（test-time adaptation, TTA）利用协议许可的当前测试输入，在产生或持续产生预测的过程中调整模型。Tent用预测熵
\[
L_{\mathrm{ent}}(\bm\theta;B)
=-\frac1{|B|}\sum_{\bm x\in B}\sum_k
p_{\bm\theta}(k\mid\bm x)\log p_{\bm\theta}(k\mid\bm x)
\]
提供无标签目标，典型实现只训练归一化层的仿射参数，并用当前批量估计归一化统计，其余权重冻结。必须区分两个时刻：原始在线流程先用当前参数得到该批输出，再反向更新，更新影响下一批；若要用更新后参数重新回答同一批，需要另做前向传播，并明确计入协议和成本。测试标签不参与这些步骤，不等于测试输入完全没有用于学习\scite{wang2021tent}

设一段无标签流为$B^{(1)},B^{(2)},\ldots$，其反馈关系是
\[
\bm\theta^{(t)}
\ \longrightarrow\ p^{(t)}(B^{(t)})
\ \longrightarrow\ L_{\mathrm{ent}}^{(t)}
\ \longrightarrow\ \bm\theta^{(t+1)}
\ \longrightarrow\ p^{(t+1)}(B^{(t+1)}).
\]
若阴影乙类已被错判为甲，降低熵可以把错误甲类概率进一步推高，下一批又以这个模型产生更偏的目标。极端地，恒预测甲具有很低的熵，却不能识别乙。逐样本更新可能缺少稳定的批量统计；批量更新有更充分统计，却会让同批样本相互影响，并引入等待和吞吐条件。持续保留状态允许积累适应，也允许积累错误；按环境重置要求环境边界可得，不能在看到结果后才决定哪些批量属于新环境。

SAR针对小批量、混合环境和类别分布变化下的失稳，筛选较可靠的样本，并通过在参数邻域考察熵损失来寻找较平坦的更新方向，减轻高敏感样本对适应的影响。其恢复机制还监测熵的运行统计，在出现塌缩迹象时恢复模型；不能把它简化为“高熵就重置”，因为退化到单类的模型可能具有过低熵。筛选阈值、可更新参数和恢复规则应预先设定，是否更准确仍需要真实任务反馈。恢复权重与清空优化器状态是两项操作，前者改变已有函数，后者改变后续学习动力学，应分别记录\scite{ket-niu2023-sar}

信息不足时，停止更新、保留最后通过验收的版本或请求少量标签都是候选措施。特别是没有办法区分输入异常和真实环境变化时，持续最小化代理目标可能使问题扩大。允许联合处理完整测试批量、逐样本在线预测和可回看的离线适应应分别报告；无标签条件不授权任意重排时间或访问未来输入。

更大范围的重估可以从头训练。记指定学习过程为$\mathcal A(S;\xi)$，$\xi$包含初始化、数据顺序及其他随机性；撤回集合为$U\subset S$，剩余数据$S^-=S\setminus U$。完整重训练以$\mathcal A(S^-;\xi')$为候选或比较参照，重建相应预处理、特征和训练状态，不能继续使用已受$U$影响却未处理的缓存。普通环境更新也可以重训练，但只有确实排除撤回记录及相关影响路径，并使用约定比较标准，才支持撤回主张。

以下排列画出一条教学训练轨迹。假设撤回记录$u$出现在第二个训练块$S_2$，$C_j$表示处理到块$j$的完整训练状态，则
\[
\begin{array}{rcl}
\text{原轨迹}&:&
C_0\xrightarrow{S_1}C_1
\xrightarrow{S_2\ni u}C_2
\xrightarrow{S_3}C_3,\\
\text{完整重训}&:&
C'_0\xrightarrow{S_1}C'_1
\xrightarrow{S_2\setminus\{u\}}C'_2
\xrightarrow{S_3}C'_3.
\end{array}
\]
从头重训重算全部阶段；若严格保证$C_1$及其数据处理和随机状态未受$u$影响，才有可能复用它来重算后缀。仅从$C_2$继续在剩余样本上训练，没有移除此前进入状态的影响。对随机算法，保持相同随机种子也未必使删除前后每步随机操作对齐，参照究竟比较一次运行还是输出分布需要另行规定。

SISA将数据分到互斥分片，各分片隔离训练，再把片内数据分成按次序处理的切片并保存快照。若$u$属于分片$a$的第二切片，删除后只需从该切片之前的片内快照重训受影响后缀，其他分片保持：
\[
\begin{array}{rcl}
\text{分片 }a&:&
C_{a,0}\to C_{a,1}
\xrightarrow{S_{a,2}\ni u}C_{a,2}\to C_{a,3},\\
\text{撤回后 }a&:&
C_{a,1}\xrightarrow{S_{a,2}\setminus\{u\}}C^-_{a,2}
\xrightarrow{S_{a,3}}C^-_{a,3},\\
\text{分片 }b&:&C_{b,3}\quad\text{保持，随后重新聚合。}
\end{array}
\]
节省来自训练依赖的隔离，代价包括多个模型、切片快照、聚合和可能不同的任务效果。其比较基准是同样分片、切片和聚合规则下排除$U$的学习过程，不能未经证明就等同于一个单体模型在全部剩余数据上重训。共享预处理若提前看过所有数据，也需要纳入依赖检查\scite{bourtoule2021sisa}

近似删除可以减少重算，但保证依赖具体模型和算法。可认证删除的一条路线在强凸正则线性模型中，用剩余目标的梯度与Hessian构造Newton修正，再控制修正后梯度残差，并结合训练时预先加入的随机扰动分析输出分布。它不是对任意神经网络做一次反向梯度就获得同样保证。若$L^-$为剩余数据目标，形式上的一步为$\bm\theta^-=\bm\theta-[\nabla^2L^-(\bm\theta)]^{-1}\nabla L^-(\bm\theta)$；是否逼近充分、噪声校准是否满足要求，需要相应假设和误差界，残差小本身不等于认证。完整机制与隐私条件见第~\ref{sec:privacy-unlearning-methods}节\scite{guo2020certifiedremoval}

选择性放弃旧能力、压低某个答案的概率和消除训练影响不能互换。模型可以给撤回照片输出错误标签，令其损失增大，却仍在中间表示中编码照片的独特背景，或保留它对其他样本预测的影响；加一条拒答规则也只改变接口响应。重新构建模型后，还要重新落实并验证其他有效编辑、行为规则和仍服务的能力。删除对象、剩余任务、重训练对照以及适用保证如何验收，由第~\ref{sec:ket-update-unlearning-validation}节展开。

\subsection{模型结构的调整}
\label{sec:ket-update-structure}

在同一组参数中反复折中，可能使新任务学不进、旧任务又守不住；现有表示容量不足时，增加训练次数也不能改变模型可表达的函数集合。结构调整改变组件数量、连接和调用的可能路径，为参数学习提供另一种配置。固定结构中的权重重置仍属于上一小节；新增模块怎样学习，也仍使用样本、损失和约束机制。

\subsubsection{组件配置与容量调整}
\label{sec:ket-update-capacity}

一种办法是将部分参数隔离。完全隔离为每个维护对象保存独立网络，旧网络不受新网络梯度影响，但存储随对象数增长，也难以自动共享新经验。部分共享让若干层共同学习，再保留专属分支，减少重复参数，却将干扰集中到共享处。共享底座加专属部分则把大部分通用计算复用，底座若冻结，新能力受现有表示限制；底座若更新，每条旧路径都可能变化。这些结构需要在同一组旧任务、相同调用协议下比较，不能把独立模型所需的任务身份视为免费信息。

Progressive Networks为新任务增加一列网络，冻结已有列，并通过横向连接让新列读取旧列的中间表示。若第$k$列第$\ell$层表示为$\bm h_k^{[\ell]}$，可概括为
\[
\bm h_k^{[\ell]}
=\sigma\!\left(
\bm W_k^{[\ell]}\bm h_k^{[\ell-1]}
+\sum_{j<k}\bm U_{k,j}^{[\ell]}\bm h_j^{[\ell-1]}
\right).
\]
方括号上标$[\ell]$在此表示网络层，不表示更新阶段；$\bm U_{k,j}^{[\ell]}$是新列读取旧表示的连接。学习新列时旧列不变，新列及横向连接接收新任务梯度；预测旧任务沿旧列，预测新任务还需计算被读取的旧列。保留发生在不变的旧调用路径上，新信息并未反向写入旧列。随着任务数增加，列数、横向连接和调用计算都可能增长，不能只统计最新列的参数\scite{ket-rusu2016-progressive}

网络增长也可以服务既有用途。例如温室与深度相机环境仍识别相同物种，却采用不同输入模块；这不必被解释为新增物种任务。新增输出或成功标准的责任仍在第~\ref{chap:ket-transfer-adaptation}章。冻结只证明被冻结参数的数值不变，如果输入预处理、共享底座、归一化统计、路由或最终校准改变，端到端行为仍可能改变。

设温室分支在温室样本上正确，新增雨季分支后，路由器把全部绿色背景照片送往雨季分支。若温室照片也有绿色背景，系统便可能退化，尽管温室分支一个参数都没有改。这个反例把两个验证对象分开：组件本身是否保持原函数，以及实际请求是否仍走到适用组件。网络有更多节点不直接证明系统获得更多可用能力。

连续增加分支不能无限回避共享取舍。假设每次新增一个同规模分支，运行若干阶段后发现多个分支在相同输入上实现近似功能，可以触发整合候选；若一个旧传感器已退役，则可以调查相关模块是否仍需保留。继续扩展增加容量和隔离，重新分配让有限容量服务新的对象，整合与压缩减少重复但重新引入函数逼近误差。保留分支并共同调用的做法见第~\ref{sec:ket-integration-composition}节，蒸馏到承接模型的做法见第~\ref{sec:ket-integration-distillation}节；本章负责确定何时触发、哪些接口必须保留以及如何验收。

回收之前应列出依赖该模块的任务、调用条件、横向连接、历史版本与外部记录。一个模块调用稀少，可能只是承担很少发生但必须识别的罕见物种，不能仅按调用率淘汰。把两个分支合并后，要检查原有输入输出接口、类别顺序和校准尺度是否兼容，旧路由引用是否更新，以及新组件是否还有后续学习容量。回收带来的存储收益应扣除替代模型、转换器和必要历史快照的成本。

\subsubsection{组件连接与路由安排}
\label{sec:ket-update-routing}

组件已经存在，仍要确定什么条件下调用哪一组。\term{路由}（routing）是从当前可见信息到组件调用的选择规则，它决定一次预测经过哪些路径。已知相机身份时可以按可信元数据选择输入模块；环境身份未知时，需要根据输入或上下文推断。任务身份说明要完成什么，环境身份说明在哪种条件下完成，预测类别则是系统要推断的答案，三者不能互相偷换。测试时将真实类别或不可见任务编号交给路由器，会改变学习问题。

对两个已保留组件$f_1,f_2$，硬路由$r(\bm x)\in\{1,2\}$产生$f_{r(\bm x)}(\bm x)$；加权调用则可在输出语义一致时写为
\[
p(y\mid\bm x)=
\alpha_1(\bm x)p_1(y\mid\bm x)
+\alpha_2(\bm x)p_2(y\mid\bm x),\qquad
\alpha_j\ge0,\quad\alpha_1+\alpha_2=1.
\]
硬路由只计算选中路径时可减少任务模块开销，但错误选择会完全绕过适用能力。加权调用可以利用多个组件，却增加推理计算，并依赖类别含义及输出尺度可比较；两个不同温度的logit不能因为维数相同就直接相加。共享底座更新会同时改变两条路径的输入，因此还要检查路由器和模块是否基于兼容的表示。

L2P在冻结的预训练模型上保存提示池，每个提示具有用于检索的键。输入先由固定的查询函数得到表示，用相似度选择若干提示，再把所选提示送入预训练模型参与任务预测。训练时利用任务损失更新提示和相关预测参数，并通过键与查询的匹配目标学习选择关系；底座冻结不表示所有学习成本为零。使用时根据输入检索，无须把真实任务身份交给系统。提示池容量、每次选择数量、键存储、查询前向和带提示前向均应计入成本；选择错误或共享提示被覆盖仍可能破坏旧能力\scite{ket-wang2021-l2p}

学习路由函数的参数属于第~\ref{sec:ket-update-parameters}节，结构层面则规定路由器能读取哪些信息、允许激活哪些组件、是否可以同时调用多条路径。输入不同而选中路径不同，不等于每次推理都重新设计了网络；既有条件计算图已经包含这些选择。图~\ref{fig:ket-update-operations}从相同初始模型出发，区分权重变化、节点连接变化与一次调用选择，防止用“模型更新”掩盖不同的代价与风险。

\begin{figure*}[htbp]
\centering
% The shared initial graph has two components A and B.
% Added nodes/edges encode capacity; line emphasis encodes one invocation.
\begin{tikzpicture}[
  x=1mm,y=1mm,font=\figurefont\small,
  box/.style={rectangle,draw=BookInk,fill=BookPaper,line width=0.7pt,
    align=center,inner sep=1.5mm,minimum height=8mm},
  module/.style={box,fill=FigureModel!10,minimum width=12mm},
  flow/.style={draw=BookInk,line width=0.7pt,-{Latex[length=1.7mm]}},
  unused/.style={flow,dashed,draw=BookMuted},
  active/.style={flow,line width=1.2pt,draw=BookTeal}
]
  \node[anchor=west,text=BookInk] at (0,122) {共同初始模型：已有组件A、B及两条可选路径};
  \node[box] (start) at (13,105) {输入};
  \node[box] (router) at (52,105) {路由};
  \node[module] (a) at (103,114) {A};
  \node[module] (b) at (103,96) {B};
  \node[box] (out) at (150,105) {输出};
  \draw[flow] (start) -- (router);
  \draw[flow] (router) -- (a);
  \draw[flow] (router) -- (b);
  \draw[flow] (a) -- (out);
  \draw[flow] (b) -- (out);
  \draw[BookRule,line width=0.6pt] (0,85) -- (163,85);

  \begin{scope}
    \node[anchor=west,font=\figurefont\small\bfseries] at (0,77)
      {A\quad 参数变化};
    \node[box] (pi) at (25,65) {输入};
    \node[box] (pr) at (25,50) {路由};
    \node[module,fill=FigureOutput!15] (pa) at (12,31) {A$^\star$};
    \node[module] (pb) at (38,31) {B};
    \node[box] (po) at (25,12) {输出};
    \draw[flow] (pi) -- (pr);
    \draw[flow] (pr) -- (pa);
    \draw[flow] (pr) -- (pb);
    \draw[flow] (pa) -- (po);
    \draw[flow] (pb) -- (po);
    \node[text=BookMuted,font=\figurefont\footnotesize] at (25,0)
      {仅A的权重改变};
  \end{scope}
  \begin{scope}[xshift=56mm]
    \node[anchor=west,font=\figurefont\small\bfseries] at (0,77)
      {B\quad 结构变化};
    \node[box] (si) at (25,65) {输入};
    \node[box] (sr) at (25,50) {路由};
    \node[box,fill=FigureModel!10] (sa) at (8,31) {A};
    \node[box,fill=FigureOutput!15] (sc) at (25,31) {C};
    \node[box,fill=FigureModel!10] (sb) at (42,31) {B};
    \node[box] (so) at (25,12) {输出};
    \draw[flow] (si) -- (sr);
    \draw[flow] (sr) -- (sa);
    \draw[flow] (sr) -- (sb);
    \draw[active] (sr) -- (sc);
    \draw[flow] (sa) -- (so);
    \draw[flow] (sb) -- (so);
    \draw[active] (sc) -- (so);
    \node[text=BookMuted,font=\figurefont\footnotesize] at (25,0)
      {新增C及其连接};
  \end{scope}
  \begin{scope}[xshift=112mm]
    \node[anchor=west,font=\figurefont\small\bfseries] at (0,77)
      {C\quad 调用变化};
    \node[box] (ri) at (25,65) {输入};
    \node[box] (rr) at (25,50) {路由};
    \node[module] (ra) at (12,31) {A};
    \node[module] (rb) at (38,31) {B};
    \node[box] (ro) at (25,12) {输出};
    \draw[active] (ri) -- (rr);
    \draw[active] (rr) -- (ra);
    \draw[unused] (rr) -- (rb);
    \draw[active] (ra) -- (ro);
    \draw[unused] (rb) -- (ro);
    \node[text=BookMuted,font=\figurefont\footnotesize] at (25,0)
      {本次选A，B未调用};
  \end{scope}
\end{tikzpicture}

\caption{参数、结构与调用变化的区别。三个面板从相同的双组件模型出发：左面板保持节点和连接，只改变组件参数；中面板增加组件及可达连接；右面板保持组件和允许的连接，以粗实线标出本次选中的路径、虚线标出未调用路径。节点增减表示结构容量变化，不以方框面积暗示参数量。三类操作可以组合，路由函数本身也可以经参数学习更新。}
\label{fig:ket-update-operations}
\end{figure*}

\subsection{外部知识的更新}
\label{sec:ket-update-external}

负责人可能频繁变更，事实还带有效时间；把每次变更都写入模型参数，并不总是最适合的维护方式。可检索记录能显式保存来源、时间和冲突，但系统必须确实调用并正确使用这些记录。外部更新因而包含内容修改与调用同步两项责任，不只是数据库中一行文字的替换。

\subsubsection{记录内容与有效范围的修改}
\label{sec:ket-update-records}

检索增强生成RAG把参数模型与外部文档库结合：检索器从库中选择与问题有关的材料，生成器在这些材料条件下回答。参数提供语言和任务处理能力，外部记录提供可替换、可追溯的内容。更换文档或索引能够改变下一次生成的条件，而不一定重新训练生成器；但原始参数中的事实仍可能与新证据竞争，生成模型也可能没有忠实使用检索材料。RAG的检索与生成形式提供一种分工，不自动解决时间有效性或数据撤回\scite{ket-lewis2020-rag}

外部编辑记忆还可以直接服务局部修正。SERAC保存编辑请求，用范围分类器判断新输入是否应受某条编辑影响：范围外调用原模型，范围内由条件生成模型利用相关编辑作答。训练需要学习哪些输入应受影响、如何依据编辑生成反事实响应；应用新请求时写入缓存，并不直接修改底模权重。与参数编辑相比，它把一部分局部性责任交给范围判断和检索，但漏选会使旧答案继续出现，过宽选择会修改无关问题，缓存增长也带来检索与冲突处理成本\scite{ket-mitchell2022-serac}

回到甲所任命示例，记录可以表示为“实体标识、关系、对象标识、有效起止时间、来源、状态”。新增一条“负责人乙，2026年7月1日起有效”，并不自动使原“负责人甲”记录过期；需要将旧记录的有效终点限定到同日。直接替换对象字段适合仅保存当前快照的副本，却会丢失历史查询依据；保存两个带时间的版本则能同时回答6月和7月的问题。标注过期表示不再用于当前时间，并非认定历史记录错误；物理删除则让该副本不再可访问，适用目的不同。

同名机构以身份标识区分，不能只按文本标题批量替换。来源的发布日期、事实生效日期和入库日期也应分开：7月3日入库的任命文件可能在7月1日已生效，7月2日入库的旧报道也不能因此覆盖任命。对相互矛盾且无法判定优先级的材料，保留各自来源和条件，将争议结论标为待决；只留下最新入库的一条，会把证据缺口伪装成确定答案。

删除记录需要沿派生关系处理副本。原文可能已生成摘要、向量、缓存答案和其他索引条目，仅删除主记录不能停止这些副本供给同一内容。若仍允许保留历史公开事实，只需改变当前适用范围；若撤回要求禁止继续使用某份记录，则应按该要求处理副本及其访问。是否还要消除它曾经造成的训练影响，必须另立参数删除目标，不能由外部记录操作隐含决定。

\subsubsection{调用结果与生成行为的同步}
\label{sec:ket-update-sync}

一次实际回答经过“问题、检索候选、适用范围判断、答案生成”。每个位置都可能使用不同版本，所以内容正确只是必要条件。解析问题时应取得主体身份和查询时间，候选检索可以同时找到旧新记录，范围判断再按有效区间及来源筛选，生成器使用适用证据而不是任意一个文本相似度最高的条目。未指定时间的“现任负责人”需要使用明确的当前时间语义，不能把模型训练时间当作今天。

任命生效后，内容操作将新旧记录关联，并更新依赖它们的“当前签批人”和“联系办公室”摘要。调用同步使这些记录的索引可检索、旧缓存失效或带上正确的时间条件，同时保留合法历史查询所需的旧记录。对7月15日查询，候选中即便含有甲，也应因时间不适用而排除；对6月15日查询，旧记录仍应被选中。负责人作为签批人的传播依据是已经给定的制度，历史审批文件上的实际签字人不能被新任命覆盖。

缓存键若只包含问题文字，不包含有效时间或版本条件，同一句“由谁负责”在任命前后的缓存就可能混用。可以按记录身份及依赖集合失效处理，也可以让缓存显式绑定所用记录版本；无论采用哪种实现，都要验证不会把旧结果误用于新条件。模型版本、索引版本和记录版本尚未全部同步时，系统可以暂停相关回答或明确其证据状态，不能让中间状态作为完整更新的验收结果。索引与缓存的工程实现归应用与第~\ref{part:systems}部分，这里约定需要保持的语义关系。

检索到两条冲突证据时，生成器应保留适用条件和不确定性，不能仅因一条较新就断言另一条完全错误。若范围分类器不能识别历史日期，或者生成器持续忽略检索证据，需要补充相应训练，并调用参数更新机制；继续向库中添加相同新事实不一定解决错误环节。验证也要分别观察候选记录、最终选中记录及答案，才能知道需要修改什么。

外部删除与参数影响消除在此仍然分开。底模可能在检索为空时继续回答已删除内容，记录不再可检索也不能说明它对其他输出的训练影响已被移除。反过来，参数重训练不会自动删除外部缓存。一次请求同时涉及两者时，应分别执行操作，并将相容的模型、记录、索引与缓存状态作为一个候选系统验收。

\section{知识更新结果的验证}
\label{sec:ket-update-validation}

完成一次训练、写入或同步操作，只得到一个候选版本。要决定是否接受，需要回到第~\ref{sec:ket-update-goals}节约定的目标：当前任务是否改善，指定修改是否在正确范围内生效，旧能力和后续学习是否维持，以及需要撤回的训练影响是否处理。它们观察的对象不同，同一更新可以通过其中几项而未通过另外几项，不能用一个总分遮盖失败。

\subsection{当前任务表现的验证}
\label{sec:ket-update-current-validation}

当前效果应与同任务、同访问条件和同调参预算下的原版本、不更新基线比较。原版本说明更新前的能力，不更新基线还应经过相同的数据解码、任务接口和评测流程，以区分模型更新与外围修复造成的变化。比较中固定目标数据的使用方式、训练和推理预算；为了选择学习率看过的数据不再充当完全独立的最终验收集，通用划分协议见第~\ref{sec:ket-evaluation-splits}节。

对花卉识别，分别报告当前雨季环境与仍需服务的温室环境，再按物种、遮挡和相机条件检查。一个教学数值例子是：原系统在温室和雨季的风险分别为$0.1,0.4$，更新后为$0.2,0.2$。等权平均由$0.25$降到$0.2$，但温室风险翻倍。若部署还要求温室风险不得增加，这个更新便未通过；若确实允许折中，也需要事先声明而不是只展示平均改善。改变两环境的评价权重还会改变汇总结论，权重必须与实际服务目标一致。

跨观测更新要另外检查新输入是否保留任务所需信息。对配对RGB和深度样本，除整体误差外，还应检查仅依赖颜色区分的类别、深度缺测区域及配准失败情形；这些局部结果可以解释为什么表示匹配成功却无法恢复原准确率。回归任务检查量纲和条件误差，时序按可用历史窗口及预测时距分组，图任务则考虑节点和邻域之间的依赖，不把高度相关样本当作相互独立的重复证据。

无标签目标只能提供代理检查。协方差差异下降、域判别准确率接近随机、熵下降或增广预测一致，都不能单独证明真实目标风险下降。没有当前真值时，可以报告这些变化、运行异常及已覆盖的假设检查，但应将真实任务验收保留为未完成，等待可靠标签或另有依据的反馈。少量标签上的改善也要报告不确定性，并考虑同植株重复照片和时间相关性；一次划分或一次随机初始化的结果不保证所有未来环境。

\subsection{修改范围与关联一致性的验证}
\label{sec:ket-update-edit-validation}

负责人变更的验收应从原请求生成测试，而不是从模型已经答对的问题反向挑选。指定表达组检查“2026年7月15日甲所负责人是谁”，预期是乙；语义改写组检查“在该日期谁主持甲所工作”，并单独加入能够核实含义一致的其他语言表达。前者失败说明目标修改尚未实现，前者成功而后者失败说明修改依赖表面提示。增加很多近乎相同的原提示，不能替代真正的改写与语言覆盖。

依赖结论组依据例子中的制度，检查7月15日项目应由谁签批、联系应转至谁的办公室，预期随负责人变为乙。无关和应保留组检查甲所成立年份、同名其他机构负责人以及6月15日甲所负责人，分别维持原有有效答案。还应检查已完成历史审批的签字人没有被新负责人覆盖。依赖组失败意味着直接事实尚未进入相关使用，应保留组失败则暴露不必要外溢或时间范围错误；局部性不是要求一切邻近问题都不变。

MQuAKE通过需要多跳事实组合的问题检查编辑之后的推理，强调回答单条已编辑事实不代表能沿事实链得到新答案\scite{ket-zhong2023-mquake}

RippleEdits进一步区分逻辑泛化、组合关系、实体别名与保留关系等传播现象\scite{ket-cohen2024-ripple}

将这些思想用于甲所例子，需要先给出有效依赖，再据此制定答案；不能把任意与机构有关的变化都解释为应当发生的“涟漪”。跨语言测试要控制翻译是否保持含义，时间测试要控制历史与当前条件，避免把语义不同的问题当成同一编辑的改写。

批量大小、连续编辑长度和事实依赖复杂度应分别变化。一批互不相关的十条编辑，与一条负责人事实连续修正十次，考验的不是同一种干扰；十条依赖组成的推理链又增加了第三种困难。可以固定每批请求数而增加阶段数，逐轮复查早先仍有效的请求；也可以固定序列长度，改变每次修改的关联范围。被后续合法请求替代的答案应按新条件验收，不能把必要覆盖计作失败，也不能只测最后一条请求。

外部更新还要观察内容、调用和响应三个位置。记录已改但命中旧缓存，说明第~\ref{sec:ket-update-sync}节的同步未完成；候选检索没有新记录，可能是索引或召回问题；检索到正确的新旧记录，却按入库时间回答历史问题，则是适用范围或生成过程有误。验收保留所用模型、库、索引和缓存版本，验证“看到了什么证据”和“用了哪项证据”。仅凭写入接口返回成功，无法定位或排除这些失败。

\subsection{能力保留与后续学习的验证}
\label{sec:ket-update-retention-validation}

持续更新需要在多个阶段复测，而不是只比较最初和最终版本。记$a_k^{(t)}$为完成阶段$t$后，在固定任务$\mathcal T_k$上按当时有效要求测得的表现，越大越好。一个按任务进入次序安排的记录可写为
\[
\begin{array}{c|ccc}
 & k=1 & k=2 & k=3\\ \hline
t=1&a_1^{(1)}&\text{--}&\text{--}\\
t=2&a_1^{(2)}&a_2^{(2)}&\text{--}\\
t=3&a_1^{(3)}&a_2^{(3)}&a_3^{(3)}
\end{array}
\]
横向表示同一版本复测哪些任务，纵向表示同一任务随更新怎样变化。破折号表示未测，不是零分；若协议允许在任务学习前做零样本评价，应另填真实观测。每项都记录样本、更新预算和测量时刻，完整评价矩阵及指标计算见第~\ref{sec:ket-evaluation-continuous}节。

保留可以比较$a_k^{(t)}$与学习完该任务时的$a_k^{(k)}$，也可以与过去最佳表现比较，但要说明参照；两种参照对应的遗忘与后向迁移量见第~\ref{sec:ket-evaluation-backward}节。若负责人在7月已经合法变更，旧测试“现任是谁”的正确答案就不能继续固定为甲；应把查询日期显式化，按有效要求重新确定参照。撤回后不再允许维持的行为也不能成为保护目标。只有在任务、时间和请求语义相容时，分数下降才可解释为相应能力退化。

可塑性要通过后续学习来检查。取原版本与更新版本的独立副本，让它们在同一新内容上使用相同样本顺序、标签数、可训练范围和更新次数，记录各预算点的损失。两个系统即使旧任务分数相同，也可能一个在少量新样本后明显改善，另一个始终学不进去；相同的保留分数无法揭示这项学习能力差异。还可定义达到预先约定风险阈值$\delta$所需的样本数$N_\delta$，但未达到时应记为在预算内未达标，不擅自外推所需样本数。

前向知识积累需要进一步比较“具有相关历史经验”和“缺少该经验”的起点。控制模型容量、当前数据、训练预算及其他预训练因素后，若前者在同样样本数下风险更低，或达到同一阈值所需样本更少，才支持历史经验帮助后续学习；这些量与学习前零样本前向迁移的区别见第~\ref{sec:ket-evaluation-forward}节。后向收益则检查学习新内容之后、没有额外给旧任务更多监督的条件下，仍有效旧任务是否改善。仅保存更多模型、增加任务数或扩张参数，不能替代这些对照。

自监督表示更新还应将两种结果分开：在更新完成后另用标签训练探测器，测得的是可读出信息；当时已有头直接产生的预测，才是即时任务能力。若两个版本的探测器训练样本或计算不同，差异也可能来自读出阶段。图~\ref{fig:ket-update-evaluation}把旧任务复测与同预算后续学习分开，前者回答保留，后者回答继续学习的能力；知识积累还需增加历史经验对照。

\begin{figure*}[htbp]
\centering
% Protocol diagram only: no invented accuracy or learning-curve data.
% Each evaluation uses a separate copy; new learning cannot alter old retests.
\begin{tikzpicture}[
  x=1mm,y=1mm,font=\figurefont\small,
  box/.style={rectangle,draw=BookInk,fill=BookPaper,line width=0.7pt,
    align=center,inner xsep=2mm,inner ysep=2mm,minimum height=12mm},
  flow/.style={draw=BookInk,line width=0.8pt,-{Latex[length=2mm]},
    rounded corners=1mm}
]
  \node[box,text width=60mm,fill=FigureModel!10] (original) at (40,93)
    {原版本\\冻结评测起点，复制两份};
  \node[box,text width=60mm,fill=FigureOutput!10] (updated) at (124,93)
    {更新版本\\冻结评测起点，复制两份};
  \node[box,text width=29mm] (old-a) at (20,62)
    {旧任务复测\\当前有效要求};
  \node[box,text width=29mm] (new-a) at (60,62)
    {同一新任务学习\\等样本、等预算};
  \node[box,text width=29mm] (old-b) at (104,62)
    {旧任务复测\\当前有效要求};
  \node[box,text width=29mm] (new-b) at (144,62)
    {同一新任务学习\\等样本、等预算};
  \draw[flow] (original.south) -- (40,79) -| (old-a.north);
  \draw[flow] (40,79) -| (new-a.north);
  \draw[flow] (updated.south) -- (124,79) -| (old-b.north);
  \draw[flow] (124,79) -| (new-b.north);
  \node[box,text width=29mm,fill=FigureInput!10] (old-ra) at (20,33)
    {原版本的\\旧任务结果};
  \node[box,text width=29mm,fill=FigureLoss!10] (new-ra) at (60,33)
    {原版本的\\后续学习过程};
  \node[box,text width=29mm,fill=FigureInput!10] (old-rb) at (104,33)
    {更新版本的\\旧任务结果};
  \node[box,text width=29mm,fill=FigureLoss!10] (new-rb) at (144,33)
    {更新版本的\\后续学习过程};
  \draw[flow] (old-a) -- (old-ra);
  \draw[flow] (new-a) -- (new-ra);
  \draw[flow] (old-b) -- (old-rb);
  \draw[flow] (new-b) -- (new-rb);
  \node[align=center,text=BookInk,text width=160mm] at (82,6)
    {两侧旧任务结果对照：能力保留；两侧同预算学习对照：可塑性\\
     进一步控制相关历史经验的有无：检验知识积累};
\end{tikzpicture}

\caption{更新前后与后续学习的对照流程。原版本和更新版本各保留独立评测副本，分别进入按当前有效要求定义的旧任务复测，以及同一新任务的学习。后续学习使用相同数据顺序、样本数和更新预算，不反过来改写旧任务复测结果。比较历史经验有无的额外对照，才能进一步判断知识积累；图中不包含假设的准确率或实测曲线。}
\label{fig:ket-update-evaluation}
\end{figure*}

\FloatBarrier
\subsection{训练影响消除的验证}
\label{sec:ket-update-unlearning-validation}

训练影响消除先固定撤回集合$U$、剩余数据$S^-$和比较学习过程$\mathcal A$。在相同任务、预处理、模型族和训练协议下生成重训练参照；若采用SISA，就使用对应的分片学习过程。撤回后仍有效的事实修改和行为要求，也应在候选与参照中按相同规则落实，否则二者差异可能只是版本请求不一致。涉及哪些随机状态和训练派生物，必须纳入这份定义，而不只是列出被删的数据文件。

删除对象上的行为、剩余任务表现和重训练接近性提供不同证据。撤回样本上的损失升高能够显示模型不再同样拟合它，但可能来自简单输出抑制；剩余任务准确说明没有观察到某类能力损害，却不说明影响已删净。比较重训练模型的预测，可以在给定输入范围内寻找残留差异，但该范围之外仍可能不同；参照本身若由随机训练产生，也需要考虑运行间变化，不能把两次随机重训的自然差异全算作删除失败。

参数接近、输出接近和训练结果分布接近不是同一个主张。参数范数距离依赖模型参数化，隐藏单元置换可以让参数很不同而函数相同；参数距离小也可能在敏感输入上改变预测。有限测试集上的输出距离只约束这些输入。更强的分布保证则比较删除算法的随机输出$Z_{\mathrm{del}}$与剩余数据学习结果$Z_{\mathrm{retrain}}$，例如在指定条件下对所有可测事件$E$要求
\[
\mathbb P(Z_{\mathrm{del}}\in E)
\le e^\varepsilon\mathbb P(Z_{\mathrm{retrain}}\in E)+\delta,
\]
并要求反方向的相应不等式。这里$\varepsilon,\delta$及两侧随机性必须来自具体保证，不能由几个提示上回答一致估出来。不同删除工作可能采用其他误差口径，应保留其原有定义，完整讨论见第~\ref{sec:privacy-unlearning-validation}节。

成员推断或其他攻击可以暴露某种残留影响，但攻击失败只表示在该攻击能力、样本和预算下未发现问题。若删除机制要求强凸目标、特定随机扰动或预先保存状态，验收时缺少这些条件，就不能沿用原保证。经验测试、机制假设与证明支持应分开记录，不能把模型拒答、某次攻击失败或缓存清空汇总为“已经认证删除”。

验收结论可以是接受、继续修正或不启用候选版本。需要硬保证却仅有经验检查时，缺口不能由提高其他任务分数抵消；观察到剩余能力损害时，应重新检查删除范围、重训练协议或候选算法。回退同样受撤回请求约束，不能返回包含已撤回影响的旧状态。表~\ref{tab:ket-update-validation}汇总这些证据的职责，用于检查是否拿错了指标。

\begin{table*}[htbp]
\small
\begin{tabularx}{\linewidth}{@{}p{19mm}XXX@{}}
\toprule
待验证主张 & 比较基线 & 主要观测对象 & 常见误判 \\
\midrule
当前效果改善 & 原版本与同协议不更新基线 & 当前及仍服务原环境的任务风险 & 代理损失下降不等于目标风险下降 \\
修改局部性 & 按有效请求确定的保留答案 & 无关实体、关系、历史时间及任务 & 参数改动少不等于语义影响小 \\
语义改写泛化 & 与请求含义一致的正确答案 & 等义改写及经核对的跨语言表达 & 原提示成功不等于不同表达均有效 \\
关联一致性 & 已核实依赖及新事实应得结论 & 有依据的依赖推理、多跳关系与版本调用 & 改写成功不等于关联推理正确 \\
能力保留 & 旧任务阶段结果或历史最佳值 & 各阶段仍有效的旧任务表现 & 旧任务不退化不等于后续仍可学习 \\
可塑性 & 同样新证据下的可比起点 & 同预算学习过程及达到阈值的代价 & 约束变弱、参数重置不等于学习恢复 \\
训练影响消除 & 剩余数据的指定学习过程 & 删除对象、剩余任务、参照距离及保证条件 & 拒答、有限输出接近、攻击失败不等于普遍删除保证 \\
\bottomrule
\end{tabularx}
\caption{更新验证中的主张与证据。各行需要独立的参照和适用条件，同一版本可能通过部分验证而未通过其他验证。知识积累还需在可塑性与保留观察之外控制历史经验差异；表中各项不合并为一个通用总分。}
\label{tab:ket-update-validation}
\end{table*}

\section{知识的持续维护}
\label{sec:ket-update-maintenance}

单次更新通过验收，仍不能保证下一轮继续有效。经验会再次变化，旧环境可能回归，请求可能被替代，记忆和模型也会占满预算。持续维护保存的不只是最新参数，还包括判断何时更新、如何恢复、哪些请求仍生效所必需的状态。具体更新调用第~\ref{sec:ket-update-parameters}至第~\ref{sec:ket-update-external}节的操作，这里关注多轮操作之间的关系。

\subsection{更新触发条件的监测}
\label{sec:ket-update-monitoring}

季节渐变时，照片亮度、湿度相关外观及类别频率可能缓慢移动；相机突然更换时，分辨率、颜色响应或深度单位可能立即跳变；短暂传感器故障则可能让若干批输入异常，设备恢复后分布回到原处。这三种情况都可能引发统计告警，却分别适合持续观察、接口复核或暂时隔离故障，不能见到变化就更新全部模型。

输入统计能说明观测发生变化，预测熵和类别频率能说明模型响应发生变化，可靠标签和任务反馈才直接支持错误关系的判断。预测类别比例突变可能是真实物种比例变化，也可能是模型塌缩；亮度变化可能无害，也可能遮掉关键纹理。监测应保留这些信号的不同解释，并与预处理错误、来源变化和标签质量一起检查。变化阈值越过时触发复核或可恢复的应对，而不是自动产生新正确标签。

标签延迟到达时，可以先记录输入窗口、版本和预测，采取停止高风险自动更新、降低处理范围或增加人工复核等措施；是否重新估计条件规律，则等待足够反馈。延迟标签应关联到产生预测时的版本和事实时间，不能用后来修订的标签口径直接解释以前的模型行为。相邻窗口高度相关时，反复告警也未必是多个独立变化。

触发策略需要比较误报、漏报和反应延迟。过敏感的阈值会频繁消耗训练与验收预算，过迟的触发则让真实失效持续更久；设置持续越界条件或多信号核对，可以改变这项取舍，但也可能延迟突发事件处理。没有明确任务边界时，窗口划分和重置条件应事先定义。上下文识别可以用于选择现有适用组件，未必需要产生新参数；监测和调度的实现细节属于第~\ref{part:systems}部分的内容范围。

\subsection{更新节奏与长期预算的安排}
\label{sec:ket-update-budget}

持续小步更新反应快，却可能在证据单一时积累偏差；积累到一批再更新能获得更好的覆盖和统计，代价是延迟；事件触发更新把计算集中在需要调整的时刻，但依赖触发信号的质量。三者应在相同经验流上记录实际处理次数、标签延迟和停机条件。所谓“在线”并不说明每条样本都只训练一次，回放和重复适应需要另外计数。

CoTTA把持续测试时适应中的误差累积纳入设计：用指数滑动平均的教师生成较稳定目标，在源模型置信度不足时利用多种增广预测的平均，训练学生匹配教师；同时随机将部分参数恢复到源模型值，以限制长期偏离。教师平均减弱单步波动，却可能随学生共同积累错误；随机恢复可以保留源能力，也可能抵消当前适应。其恢复基准若早于一项有效事实修正或数据撤回，还可能重新引入禁止状态，因此在带请求维护中必须更换为合规恢复基准或禁用不相容恢复\scite{wang2022cotta}

回归温室时，应检查雨季更新后的系统能否重新使用温室能力，以及恢复需要多少样本、更新和时延。保留旧模型可以减少重新学习，但增加存储；只保留小记忆节省模型空间，却需要重新优化；压缩记忆可能丢失罕见情形。对历史记录还要持续入库、去重、补齐延迟标签和处理有效期，不能让过时负责人答案继续作为无条件回放目标，也不能将已撤回样本重新选入记忆。

长期存储账目可以写为
\[
B_{\mathrm{total}}
=B_{\mathrm{model}}+B_{\mathrm{memory}}+B_{\mathrm{state}}
+B_{\mathrm{external}}+B_{\mathrm{recovery}}.
\]
这五项分别覆盖模型与路由、样本及历史输出、优化器与统计、外部记录及索引、恢复快照；共享对象只计一次并注明归属。作为教学账目，设五项依次占$400,250,100,100,150$个统一存储单位，总额1000。若新增分支需要80单位，其他项不变且总额固定，历史记忆就只能从250降到170；若不接受这项覆盖损失，需要压缩、回收其他对象或拒绝扩容，不能仍声称保持原记忆预算。

按上述五项的固定顺序，这次调整的账目为
\[
(400,250,100,100,150)
\ \longrightarrow\
(480,170,100,100,150),
\qquad B_{\mathrm{total}}=1000.
\]
模型多占的80单位来自记忆的80单位让渡；新增材料仍要在剩余170单位内通过保留和替换进入，不能另开一个未计费的记忆池。

新材料入库还会发生“新增记录、保留比较、替换旧记录”的循环。替换应同时处理其标签、特征和缓存预测，延迟标签到达后再更新对应记录；若只按照片数量计算，忽略随类别增多而变宽的logit缓存，字节预算仍会增长。固定每类内存、固定总内存和固定模型数量是不同条件。推理的活跃组件数、检索时延、增广次数、更新频率和恢复重算也要形成长期账目，不能只给单轮训练时间。

\subsection{知识版本与有效请求的跟踪}
\label{sec:ket-update-versions}

版本记录要能回答“这个响应为什么在这个时间成立”。除模型差异外，还需保存请求依据、主体身份、适用区间、被改对象、依赖结论、验证结果和未决问题。某次验收只对当时的请求集合和组件版本有效；新请求改变了其中一个条件，就需要标记哪些旧结论待复查，而不是把过去的通过状态永久沿用。

在教学例子中，第一次任命使乙从7月1日起负责，第二次任命使丙从9月1日起负责。第二次更新应把乙的适用区间收束为$[\text{7月1日},\text{9月1日})$，并新增丙的区间；它不应使7月15日的历史回答变成丙。若后来收到另一来源声称9月应由丁代理，需要核对机构范围、代理权限和生效时间，不能只按请求先后覆盖。冲突尚未解决时保留分歧和待决状态，依赖该事实的签批结论也相应受限。

参数编辑还可能发生非语义性的顺序干扰。第二次在共享层写入内容时，第一次编辑依赖的键或下游表示可能改变，即使两条请求指向不同实体也未必互不影响。因此连续维护按阶段复测仍有效请求和相关保留集，并记录批量规模、序列长度及依赖复杂度。一次秩一修改的矩阵增量可以保存，但以后网络已继续训练，直接减去原增量并不保证函数回到当时状态；可逆范围要由实际依赖与验证确定。

一次响应可标识为版本元组
\[
\begin{aligned}
v=(&v_{\mathrm{data}},v_{\mathrm{prep}},v_{\mathrm{embed}},
v_{\mathrm{train}},v_{\mathrm{model}},v_{\mathrm{eval}},\\
&v_{\mathrm{route}},v_{\mathrm{deploy}},v_{\mathrm{record}},
v_{\mathrm{index}},v_{\mathrm{cache}}),
\end{aligned}
\]
其中前八项依次记录数据快照、预处理、嵌入或特征生成、训练配置、模型、评价协议、路由和部署包，后三项记录外部事实链路。元组还关联当前有效请求集合$\mathcal E_{\mathrm{active}}$和查询时间。模型相同而索引不同，可能找到不同证据；数据相同而预处理或嵌入不同，也可能形成不同模型与检索结果。图~\ref{fig:ket-update-versions}用独立的模型与记录节点表达这些依赖。只保留最新模型文件，无法重建当时使用的历史材料、评价条件、调用配置和仍需履行的请求。

\begin{figure*}[htbp]
\centering
% Teaching versions: e is an active fact edit; u is a withdrawn record.
% Dotted rollback path is a candidate only, never a direct deployment edge.
\begin{tikzpicture}[
  x=1mm,y=1mm,font=\figurefont\small,
  box/.style={rectangle,draw=BookInk,fill=BookPaper,line width=0.7pt,
    align=center,inner xsep=2mm,inner ysep=2mm,minimum height=13mm},
  flow/.style={draw=BookInk,line width=0.8pt,-{Latex[length=2mm]},
    rounded corners=1mm},
  dependency/.style={flow,dashed,draw=BookBlue},
  candidate/.style={flow,densely dotted,draw=BookAmber,line width=1pt}
]
  \node[anchor=west,text=BookInk] at (6,118) {模型版本：按时间更新};
  \node[box,text width=29mm,fill=FigureModel!10] (m0) at (24,101)
    {$M^{(0)}$\\未满足$e,u$};
  \node[box,text width=29mm,fill=FigureModel!10] (m1) at (82,101)
    {$M^{(1)}$\\已落实$e$};
  \node[box,text width=29mm,fill=FigureOutput!10] (m2) at (140,101)
    {$M^{(2)}$\\已处理$u$};
  \draw[flow] (m0) -- node[above,font=\figurefont\footnotesize] {修正$e$} (m1);
  \draw[flow] (m1) -- node[above,font=\figurefont\footnotesize] {撤回$u$} (m2);

  \node[anchor=west,text=BookInk] at (6,81) {外部记录版本：保留有效历史};
  \node[box,text width=29mm] (d0) at (24,65)
    {$D^{(0)}$\\旧记录范围};
  \node[box,text width=29mm] (d1) at (82,65)
    {$D^{(1)}$\\新旧时间区间};
  \node[box,text width=29mm] (d2) at (140,65)
    {$D^{(2)}$\\处理撤回副本};
  \draw[flow] (d0) -- node[above,font=\figurefont\footnotesize] {$e$同步} (d1);
  \draw[flow] (d1) -- node[above,font=\figurefont\footnotesize] {$u$同步} (d2);
  \node[box,text width=48mm,fill=FigureLoss!10] (check) at (55,18)
    {重新核对全部有效请求\\补做$e$及$u$的合规处理};
  \node[box,text width=42mm,fill=FigureOutput!10] (restored) at (136,18)
    {恢复候选$M^{\mathrm R}$\\验收通过后才能启用};
  \draw[candidate] (m0.west) -- (2,101) -- (2,18) -- (check.west);
  \node[anchor=west,text=BookAmber,font=\figurefont\footnotesize]
    at (8,40) {旧快照只作候选来源};
  \draw[flow] (check) -- (restored);
  \draw[dependency] (d2.south) -- (140,42) -| (restored.north);
  \node[anchor=east,text=BookBlue,font=\figurefont\footnotesize]
    at (136,42) {当前记录相容};
  \node[align=center,text=BookMuted,font=\figurefont\footnotesize,
    text width=156mm] at (82,-3)
    {实线：版本更新；虚线：相容依赖；点线：待审核恢复路径};
\end{tikzpicture}

\caption{连续修改与受约束回退。模型和外部记录分别沿实线形成版本，虚线表示请求或版本依赖，点线表示待审核的回退候选。教学示例中的$e$为仍有效的事实修改，$u$为撤回记录；旧模型虽可能恢复某项性能，却不能绕过$e$与$u$直接启用。恢复候选须与当前适用的外部记录相容，并重新验证全部有效请求；图不暗示外部记录删除等同于参数影响消除。}
\label{fig:ket-update-versions}
\end{figure*}

\FloatBarrier
\subsection{异常恢复与回退结果的复核}
\label{sec:ket-update-rollback}

回退是另一项系统变更，而不是跳出维护要求的例外。先确定需要恢复哪项能力以及哪些快照可用，再列出在恢复时仍有效的事实编辑、行为约束和撤回集合。候选快照如果早于这些请求，必须重新应用尚未满足的要求，或者从合规材料重建；无法保证合规时，不能直接启用。权重、路由、外部库、索引和缓存都需处于相容状态。

假设模型$M^{(0)}$在温室表现较好，但含有后来撤回的记录$u$的训练影响；$M^{(1)}$修正负责人事实，$M^{(2)}$又完成对$u$的合规处理。一次后续适应使当前性能下降，直接恢复$M^{(0)}$虽可能恢复温室效果，却同时恢复旧负责人答案和被撤回影响。正确候选应由合规快照出发，或对旧快照重新执行所需事实修正与删除过程，得到$M^{\mathrm R}$，再进行完整验收。删除算法若依赖旧训练状态而状态已不可用，就不能承诺可以低成本补做。

外部库也不能任意跟随模型回到旧时间。历史负责人记录可以保留服务历史查询，但当前记录范围仍应满足有效任命；已撤回的缓存或索引内容不能因回滚备份重新开放。模型采用新的有效请求、外部库却回到旧版本，会形成新的不一致。恢复测试应同时覆盖目标能力、历史语义、无关响应、撤回对照及版本组合，接受之后继续观察回归环境和后续学习。

停止自动更新也需要明确条件。连续代理损失改善而可靠任务反馈恶化、请求来源冲突未解决、没有可验证的合规恢复路径，或长期预算不足以同时保留必需能力时，可以暂停相关更新，限制服务范围并等待新证据。应记录哪些能力仍可提供、哪些问题待决和何时重新评估，不把暂停描述为已经完成修复。日志、存储和部署的具体机制由系统卷实现，知识层面的责任是让当前行为始终对应到可说明的版本、条件和证据。

\FloatBarrier
\section{本章小结}
\label{sec:ket-update-summary}

雨季花卉识别失效可能来自频率、条件规律或观测信息变化，需要不同证据才能确定更新目标；负责人变更则要求明确实体、时间和依赖，训练记录撤回又提出不同于改写答案的影响消除目标。它们可以同时进入一次更新，但不能由相同的操作完成状态或单项分数共同验收。

样本取用、学习目标、可变参数和计算过程共同形成参数更新；结构调整改变可用组件与调用关系，外部更新改变可检索内容及适用范围。回放与行为约束可以互补，冻结参数不能保护错误路由，外部记录正确也不能保证生成答案正确。保留要求必须服从有效修正和撤回，局部性限制不必要外溢，同时允许有依据的关联结论改变。

持续维护还要分别检验旧能力、当前收益和后续学习，并以可比对照判断知识是否积累。快照、记忆和外部库只有与当前有效请求一起管理，才能支持受约束的恢复；旧版本分数较好不构成重新引入已撤回影响的理由。用途扩展时由第~\ref{chap:ket-transfer-adaptation}章确定新任务，扩展后的能力保留、版本相容与长期成本则继续接受这里的维护要求。
